\documentclass[aps,prd,twocolumn,superscriptaddress,nofootinbib]{revtex4-2}
\usepackage{amsmath}
\usepackage{amssymb}
\usepackage{amsthm}
\usepackage{booktabs}
\usepackage{xcolor}
\usepackage{longtable}
\usepackage{multirow}
\usepackage{array}
\usepackage{float}
\usepackage{graphicx}
\graphicspath{{../figures/}}

\newtheorem{definition}{Definition}
\newtheorem{proposition}{Proposition}

\usepackage{hyperref}
\hypersetup{
    pdftitle={Spieltheoretische Kosmologie und das Curvature Relaxation Model},
    pdfauthor={Lukas Geiger},
    colorlinks=true,
    linkcolor=black,
    urlcolor=blue,
    citecolor=black
}

\begin{document}

% ===================================================================
% TITELSEITE
% ===================================================================

\title{Spieltheoretische Kosmologie und das Krümmungs-Rückgabemodell: Nash-Gleichgewichte zwischen Nullraum und Raumzeitblase als Erklärungsrahmen für die beschleunigte Expansion}

\author{Lukas Geiger}
\email{lukisch@users.noreply.github.com}
\affiliation{Independent Researcher, Bernau, Germany}

\date{\today}


\begin{abstract}
Wir entwickeln einen spieltheoretischen Rahmen für die Kosmologie, in dem die Raumzeitentwicklung als Nash-Gleichgewicht zwischen einem metastabilen Quantenvakuum (Nullraum) und einer Raumzeitblase modelliert wird. Das zentrale Ergebnis ist das \textit{Curvature Relaxation Model} (CRM), das die beschleunigte Expansion nicht durch Dunkle Energie erklärt, sondern durch ein nachlassendes Krümmungs-Rückgabepotential $\Omega_\Phi(a) = \Phi_0 \cdot \tanh(k(a - a_{\mathrm{trans}}))$ -- ein geometrisches ``Gedächtnis'' der Urknall-Energiekonzentration. Getestet gegen 1.590 Pantheon+ SN-Ia-Datenpunkte \cite{Scolnic2022} mit diagonalen und vollen Kovarianzfehlern liefert das CRM $\Delta\chi^2 = -12{,}2$ ($-11{,}2$) und $\Delta\mathrm{AIC} = -8{,}2$ ($-7{,}2$) gegenüber $\Lambda$CDM; der 5-Fold-Kreuzvalidierungs-Mittelwert liegt leicht niedriger, ist ohne Fold-Varianzen aber als vergleichbar statt entscheidend zu werten. Die MCMC-Analyse ergibt $\Omega_m = 0{,}368 \pm 0{,}024$; vier alternative Sättigungsfunktionen liefern vergleichbare Ergebnisse und belegen die Robustheit. Der phänomenologische Pantheon+-Fit erzeugt eine phantomartige Effektivdiagnostik, $w_{\mathrm{eff}}(z)<-1$, aber aktuelle DESI-DR2-Kombinationen bevorzugen $w_0>-1$, $w_a<0$ \cite{DESI2025}; persistente Phantomhaftigkeit wird daher als aufzulösende Spannung behandelt, nicht als fundamentale CRM-Vorhersage. Die robustere Spätsignatur ist ein späterer Beschleunigungsübergang ($z_{\mathrm{acc}} = 0{,}52$ vs.\ $0{,}84$), testbar mit Euclid, Roman und DESI. Die spätere Beschleunigung impliziert eine verlängerte materiedominierte Wachstumsphase, die die JWST ``Universe Breakers'' bei $z > 7$ \cite{Labbe2023} und die El-Gordo-Anomalie \cite{Asencio2023} natürlich erklärt. Der beschränkte Fit-Hintergrund nähert sich ohne Big Rip einem de-Sitter-Zustand; Störungsstabilität ist durch dieses Sättigungsargument nicht nachgewiesen. Die spieltheoretischen Axiome besitzen eine exakte thermodynamische Äquivalenz in der Jacobson-Tradition \cite{Jacobson1995}. Der Analysecode ist öffentlich verfügbar.\footnote{\url{https://github.com/research-line/crm-cosmology}}

\end{abstract}

\keywords{Spieltheorie, Nash-Gleichgewicht, Kosmologie, Dunkle Energie, Krümmungs-Rückgabepotential, Curvature Relaxation Model, Friedmann-Gleichung, Finsler-Gravitation, beschleunigte Expansion, Zustandsgleichung}

\thanks{KI-Werkzeuge (Claude, Anthropic; Copilot, Microsoft/GitHub; Gemini, Google DeepMind) wurden für mathematische Formalisierung, Code-Entwicklung und Texterstellung verwendet. Alle physikalischen Hypothesen, wissenschaftliche Interpretation und Verantwortung für die Inhalte liegen vollständig bei dem Autor. Der Analysecode ist verfügbar unter \url{https://github.com/research-line/crm-cosmology}.}

\thanks{Arbeit~I in der CRM-Serie.}

\maketitle


% ===================================================================
% 1. EINLEITUNG
% ===================================================================
\section{Einleitung}
\label{sec:introduction}

Die Entdeckung der beschleunigten Expansion des Universums durch die Beobachtung entfernter Typ-Ia-Supernovae im Jahr 1998 durch die Teams um Perlmutter \cite{Perlmutter1999} sowie Riess und Schmidt \cite{Riess1998} markiert einen Wendepunkt der modernen Kosmologie. Für diese Entdeckung wurde 2011 der Nobelpreis für Physik verliehen. Das Standardmodell der Kosmologie, $\Lambda$CDM, erklärt die Beschleunigung durch eine kosmologische Konstante~$\Lambda$, die etwa 68\,\% der Energiedichte des Universums ausmacht \cite{Planck2020}. Trotz seiner empirischen Erfolge steht $\Lambda$CDM vor tiefgreifenden konzeptuellen Problemen:

\begin{enumerate}
\item \textbf{Das Kosmologische-Konstante-Problem:} Die beobachtete Vakuumenergiedichte ist um $\sim$60--120 Größenordnungen kleiner als theoretische Vorhersagen der Quantenfeldtheorie \cite{Weinberg1989}.
\item \textbf{Das Koinzidenz-Problem:} Warum sind $\Omega_m$ und $\Omega_\Lambda$ gerade in der heutigen Epoche von vergleichbarer Größenordnung?
\item \textbf{Die $H_0$-Spannung:} Die lokale Messung des Hubble-Parameters ($H_0 \approx 73$\,km/s/Mpc) weicht signifikant von der aus CMB-Daten abgeleiteten ($H_0 \approx 67{,}4$\,km/s/Mpc) ab \cite{Planck2020, Riess2022}.
\end{enumerate}

Jüngste Resultate des \textit{Dark Energy Spectroscopic Instrument} (DESI) verstärken die Zweifel an einer strikt konstanten Dunklen Energie: Die DR1-Analyse baryonischer akustischer Oszillationen in Kombination mit CMB- und Supernova-Daten zeigt eine Präferenz von 2,5--3,9$\sigma$ für ein Modell mit zeitabhängigem Zustandsgleichungsparameter $w(z)$ gegenüber $\Lambda$CDM \cite{DESI2024}. Das nachfolgende DESI-DR2-Resultat \cite{DESI2025} verstärkt diesen Befund: Im $w_0w_a$CDM-Modell bevorzugt es den Quadranten $w_0>-1$, $w_a<0$, mit supernovaabhängigen Präferenzen von 2,8--4,2$\sigma$ gegenüber $\Lambda$CDM. Beispielsweise liefert DESI+CMB+Union3 $w_0=-0{,}667\pm0{,}088$ und $w_a=-1{,}09^{+0{,}31}_{-0{,}27}$. Das steht in möglicher Spannung zur phantomartigen Effektivdiagnostik des phänomenologischen CRM-Fits ($w_\mathrm{eff}< -1$); die Einschränkungen werden in Abschnitt~\ref{sec:discussion} diskutiert.

Parallel dazu zeigen theoretische Arbeiten, dass die Beschleunigung auch ohne Dunkle Energie erklärbar sein könnte: Pfeifer et al.\ \cite{Pfeifer2025} demonstrieren im Rahmen der Finsler-Gravitation, dass eine verallgemeinerte Raumzeitgeometrie natürlicherweise eine exponentielle Expansion im Vakuum erzeugt. Trivedi und Venkatasubramanian \cite{Trivedi2025} zeigen in ihrer \textit{Cosmological Teleodynamics}, dass das Universum wie ein ``riesiges Potentialspiel'' operiert und sich einem kontinuierlichen Nash-Gleichgewicht annähert, wobei die kosmische Beschleunigung als emergenter Effekt dynamischen Gedächtnisses in einem selbstgravitierenden Medium erscheint.

Die vorliegende Arbeit verknüpft diese Entwicklungen mit einem eigenständigen Ansatz: Ausgehend von einer spieltheoretischen Modellierung der Wechselwirkung zwischen Quantenvakuum und Raumzeit wird das \textit{Curvature Relaxation Model} (CRM; dt.\ Krümmungs-Rückgabemodell) entwickelt, das die beschleunigte Expansion als ``nachlassende Bremse'' statt als ``neuen Antrieb'' interpretiert.

% ===================================================================
% 2. SPIELTHEORETISCHER RAHMEN
% ===================================================================
\section{Spieltheoretischer Rahmen: Nullraum und Raumzeitblase}
\label{sec:gametheory}

\textit{Epistemologischer Status.} Der in diesem Abschnitt entwickelte spieltheoretische Rahmen ist eine \emph{heuristisch-konzeptionelle} Interpretationsschicht in der Tradition von Jacobsons thermodynamischer Ableitung der Einstein-Gleichungen \cite{Jacobson1995}. Er liefert physikalische Intuition und motiviert die funktionale Form des CRM, ist aber nicht der operative dynamische Rahmen. Die operative Dynamik ist die Lagrange-Formulierung $\mathcal{L} = R/(16\pi G) + \gamma R^2 + \frac{1}{2}(\partial\phi)^2 - V_{\mathrm{PT}}(\phi)$, die in Paper~III \cite{Geiger2026c} abgeleitet wird und aus der alle testbaren Vorhersagen folgen. Das spieltheoretische Bild sollte daher als \emph{interpretatives Meta-Framework} verstanden werden, nicht als alternative fundamentale Dynamik.

\subsection{Grundannahmen}
\label{subsec:assumptions}

Der hier vorgeschlagene Rahmen geht von folgenden Annahmen aus:

\begin{enumerate}
\item Es existiert ein metastabiler Quantenvakuumzustand (im Folgenden: \textit{Nullraum}), der durch Quantenfluktuationen charakterisiert ist; beobachtbare Vakuumeffekte wie die Casimir-Kraft dienen dabei als minimale physikalische Analogie \cite{Casimir1948}. \textit{Ontologische Einordnung:} Der Nullraum ist eine \emph{prä-geometrische Abstraktion} -- eine Randbedingung für das Variationsproblem -- kein konkretes physikalisches Objekt mit eigenständiger raumzeitlicher Existenz. Mehrere UV-Completion-Programme (Schleifen-Quantengravitation, Causal-Set-Theorie, Quanten-Fehlerkorrektur; siehe Paper~III \cite{Geiger2026c}, Anhang) sind mit dieser Abstraktion vereinbar.
\item Eine außergewöhnlich große Fluktuation entnimmt dem Nullraum einmalig eine Energiemenge~$E_0$, die einen Konzentrationsgradienten erzeugt.
\item Zur Einkapselung und kontrollierten Neutralisation dieses Gradienten entsteht Raumzeit als dynamische Struktur -- die \textit{Raumzeitblase} (Tochtersystem).
\item Zwischen Nullraum (Muttersystem) und Raumzeitblase besteht ein spieltheoretisches Gleichgewicht.
\end{enumerate}

Diese Annahmen werden im Folgenden in einen formalen Rahmen überführt.

\subsection{Akteure und Ziele}
\label{subsec:agents}

Das System wird als Zweipersonen-Potentialspiel modelliert:

\begin{description}
\item[\textbf{Nullraum (Muttersystem):}] Primärziel ist der Selbstschutz -- die Erhaltung seiner strukturellen Integrität. Er reguliert die Kopplungsstärke zur Raumzeitblase über effektive Randbedingungen (``Gatekeeping''), erzwingt langsame Energieabfuhr (Dämpfung) und bildet Pufferzonen (horizontartige Hüllen).
\item[\textbf{Raumzeitblase (Tochtersystem):}] Primärziel ist die kontrollierte Rückkehr in den Nullzustand bei gleichzeitigem Schutz des Muttersystems. Die Strategien umfassen kaskadierten Gradientenabbau, adiabatische Rückführung und Entropiemanagement.
\end{description}

\subsection{Mathematische Formulierung als Potentialspiel}
\label{subsec:potentialgame}

Die globale Zielfunktion des Systems lautet:
\begin{equation}
\Phi = \alpha \cdot S_{\mathrm{Mutter}} + \beta \cdot R_{\mathrm{Tochter}} - \gamma \cdot G
\label{eq:potential}
\end{equation}
wobei $S_{\mathrm{Mutter}}$ die strukturelle Integrität des Nullraums, $R_{\mathrm{Tochter}}$ den Rückkehrfortschritt und $G$ den verbleibenden Konzentrationsgradienten beschreibt; $\alpha, \beta, \gamma > 0$.

\begin{definition}[Nash-Gleichgewicht des kosmologischen Spiels]
Ein Strategiepaar $(s_M^*, s_T^*)$ von Nullraum und Raumzeitblase bildet ein Nash-Gleichgewicht, wenn gilt:
\begin{align}
\Phi(s_M^*, s_T^*) &\geq \Phi(s_M, s_T^*) \quad \forall\, s_M \\
\Phi(s_M^*, s_T^*) &\geq \Phi(s_M^*, s_T) \quad \forall\, s_T
\end{align}
Keine Seite kann durch einseitige Abweichung von ihrer Strategie das Gesamtpotential verbessern, ohne die Stabilität des Systems zu gefährden.
\end{definition}

Dies ist das von Nash eingeführte Kriterium einseitiger Nicht-Verbesserbarkeit \cite{Nash1950}, das hier nur als interpretative Vorlage eines Zwei-Spieler-Potentialspiels verwendet wird.

Der zentrale \textbf{Zielkonflikt} besteht darin, dass eine zu schnelle Reduktion von~$G$ (sofortige Rückkehr) $S_{\mathrm{Mutter}}$ gefährdet, während eine zu langsame Reduktion die Entropie und die Kosten innerhalb der Blase erhöht. Das Nash-Gleichgewicht erzwingt daher eine kontrollierte, zeitlich gestreckte Neutralisation.

\subsection{Emergente Gesetze aus dem Gleichgewicht}
\label{subsec:emergentlaws}

\textit{Vorbehalt zum spieltheoretischen Rahmenwerk:} Die in den Abschnitten~\ref{subsec:assumptions}--\ref{subsec:thermodynamic} vorgestellte spieltheoretische Konstruktion liefert eine \emph{interpretative Motivation} für die CRM-Phänomenologie -- sie stellt keine formale Herleitung des CRM-Potentials $\Omega_\Phi(a)$ aus ersten Prinzipien dar. Die funktionale Form $\Omega_\Phi(a) \propto \tanh(k(a - a_\mathrm{trans}))$ wurde gewählt, um Typ-Ia-Supernova-Daten anzupassen; das spieltheoretische Narrativ wurde konstruiert, um diese Wahl physikalisch intuitiv zu machen. Dies ist ein legitimer methodischer Ansatz in der phänomenologischen Kosmologie, Leser sollten sich jedoch bewusst sein, dass das spieltheoretische Argument die funktionale CRM-Form nicht eindeutig festlegt. Eine formale Herleitung aus einer Quantengravitations- oder effektiven Feldtheorie-Wirkung steht noch aus.

Aus der spieltheoretischen Gleichgewichtsbedingung emergieren physikalische Gesetzmäßigkeiten (als heuristische Konsequenzen des Nash-Gleichgewichtsbildes):

\begin{enumerate}
\item \textbf{Energieerhaltung:} Konservative Feldgleichungen entstehen als notwendige Bedingung für stabilen Gradientenabbau.
\item \textbf{Kausalstruktur:} Die Hüllenbildung des Nullraums erzwingt eine maximale Ausbreitungsgeschwindigkeit für Informationen und Energie.
\item \textbf{Entropischer Zeitpfeil:} Die ``Zeit'' innerhalb der Blase ist die Ordnung, entlang der der Konzentrationsgradient nivelliert wird.
\item \textbf{Flusslimitierung:} Maximalflüsse über die Hülle skalieren sublinear mit dem internen Überschuss und verhindern Runaway-Prozesse.
\item \textbf{Asymptotische Rückkehr:} Der Restgradient $G \to 0$ nähert sich nur asymptotisch; es gibt kein katastrophales Finale.
\end{enumerate}

Die letzte Eigenschaft ist für die Kosmologie besonders bedeutsam: Sie impliziert, dass die Expansion des Universums sich nie umkehrt, sondern asymptotisch abläuft -- konsistent mit den beobachteten Daten.

\subsection{Thermodynamische Äquivalenz}
\label{subsec:thermodynamic}

Der spieltheoretische Rahmen der vorangegangenen Abschnitte besitzt eine präzise Äquivalenz in der Sprache der Thermodynamik und Variationsphysik. Diese Verbindung spiegelt einen tiefen mathematischen Isomorphismus zwischen Spieltheorie und statistischer Mechanik wider \cite{Wolpert2006}:

\begin{itemize}
\item \textbf{Nash-Gleichgewicht} $\leftrightarrow$ \textbf{Thermodynamisches Gleichgewicht:} Der Zustand, von dem aus kein Teilsystem seine Zielfunktion einseitig verbessern kann, entspricht dem Zustand maximaler Entropie unter Zwangsbedingungen.
\item \textbf{Potentialfunktion~$\Phi$} $\leftrightarrow$ \textbf{Entropiefunktional:} Die globale Größe, die durch die Gleichgewichtsstrategien extremiert wird, ist die Entropie des gekoppelten Systems.
\item \textbf{Optimale Strategiewahl} $\leftrightarrow$ \textbf{Prinzip der kleinsten Wirkung:} Die Strategie jedes Akteurs entspricht Feldgleichungen aus $\delta S = 0$.
\item \textbf{Gradientenabbaurate} $\leftrightarrow$ \textbf{Prinzip maximaler Entropieproduktion (MEPP):} Unter allen mit den Randbedingungen verträglichen Pfaden entwickelt sich das System entlang desjenigen, der die Entropieproduktion maximiert \cite{Dewar2003}.
\end{itemize}

Diese Äquivalenz gründet auf der Thermodynamik schwarzer Löcher \cite{Hawking1974} und auf einem fundamentalen Resultat: Jacobson \cite{Jacobson1995} zeigte, dass die Einsteinschen Feldgleichungen aus der Proportionalität von Entropie und Horizontfläche zusammen mit der Clausius-Relation $\delta Q = T\,dS$ abgeleitet werden können, was die Allgemeine Relativitätstheorie selbst als thermodynamische Zustandsgleichung etabliert. Neuere Arbeiten haben dieses Paradigma auf nicht-Riemannsche Geometrien mit Torsion und Nicht-Metrizität erweitert~\cite{JacobsonNonRiemannian2026}, was darauf hindeutet, dass der thermodynamische Zugang zur Gravitation allgemeiner ist als der metrische Rahmen und die zusätzlichen geometrischen Freiheitsgrade (wie den CRM-Scalaron) als entropie-maximierende Kanäle natürlich berücksichtigen kann. Wenn die Standard-Gravitation aus der Thermodynamik folgt, ist es natürlich zu erwarten, dass \textit{modifizierte} Gravitation -- hier der Krümmungs-Rückgabemechanismus -- aus modifizierten thermodynamischen Randbedingungen folgt, nämlich der Zwangsbedingung, dass das Muttervakuum (Nullraum) strukturell stabil bleibt, während die Entropieproduktion im Tochtersystem maximiert wird.

Der Sättigungsmechanismus aus Abschnitt~\ref{subsec:saturation} ist damit direkt als Annäherung an das thermodynamische Gleichgewicht interpretierbar: Die $\tanh$-Dynamik beschreibt die Entropieproduktionsrate eines Systems mit endlicher Kapazität, das sich unter der gemeinsamen Zwangsbedingung des Zweiten Hauptsatzes und der Muttersystem-Stabilität entwickelt.

% ===================================================================
% 3. DAS CURVATURE RELAXATION MODEL
% ===================================================================
\section{Das Curvature Relaxation Model (CRM)}
\label{sec:crm}

\subsection{Physikalische Motivation}
\label{subsec:motivation}

Im spieltheoretischen Rahmen der vorigen Sektion wird die Raumzeit als ``Bremsmechanismus'' interpretiert, der die sofortige Rückkehr der Energie in den Nullraum verhindert. Die zentrale physikalische Einsicht lautet:

\begin{quote}
\textit{Die beobachtete beschleunigte Expansion ist nicht durch eine neue Energieform verursacht, sondern durch ein nachlassendes Krümmungs-Rückgabepotential -- eine Art geometrisches ``Gedächtnis'' der anfänglichen Energiekonzentration beim Urknall.}
\end{quote}

Die Analogie ist die einer gespannten Feder: Anfänglich herrscht maximale Spannung (hohe Krümmung) mit starker Rückstellkraft. Mit der Zeit lässt die Spannung nach, die Rückstellkraft nimmt ab, und die Expansion ``beschleunigt'' relativ zur gebremsten Frühphase -- wie ein Auto, bei dem die Handbremse langsam gelöst wird.

\subsection{Modifizierte Friedmann-Gleichung}
\label{subsec:friedmann}

Die Standard-Friedmann-Gleichung im $\Lambda$CDM-Modell lautet:
\begin{equation}
H^2(a) = H_0^2 \left[\Omega_m\,a^{-3} + \Omega_\Lambda\right]
\label{eq:friedmann_lcdm}
\end{equation}

Im CRM wird die kosmologische Konstante durch ein zeitabhängiges Krümmungs-Rückgabepotential ersetzt:
\begin{equation}
H^2(a) = H_0^2 \left[\Omega_m\,a^{-3} + \Omega_\Phi(a)\right]
\label{eq:friedmann_cfm}
\end{equation}

Das Krümmungs-Rückgabepotential ist definiert als:
\begin{equation}
\Omega_\Phi(a) = \Phi_0 \cdot \frac{\tanh\!\big(k\cdot(a - a_{\mathrm{trans}})\big) + s}{1 + s}
\label{eq:potential_cfm}
\end{equation}
wobei $s = \tanh(k \cdot a_{\mathrm{trans}})$ ein Normierungsshift ist, der $\Omega_\Phi(0) = 0$ sicherstellt, und:
\begin{itemize}
\item $a$ der Skalenfaktor ist ($a=1$ heute, $a \to 0$ beim Urknall),
\item $\Phi_0$ die Amplitude (aus der Flachheitsbedingung $\Omega_m + \Omega_\Phi(1) = 1$ abgeleitet),
\item $k$ die Übergangsschärfe,
\item $a_{\mathrm{trans}}$ der Übergangsskalenfaktor.
\end{itemize}
Die konkreten Parameterwerte werden in Abschnitt~\ref{sec:numerics} aus dem Fit gegen den Pantheon+-Datensatz bestimmt.

\subsection{Dynamischer Sättigungsmechanismus}
\label{subsec:saturation}

Die $\tanh$-Parametrisierung ist keine \textit{ad hoc} gewählte Fitfunktion, sondern entsteht als exakte Lösung eines physikalisch motivierten \textbf{dynamischen Sättigungsmechanismus}. Die zentrale Annahme lautet: Die Raumzeitblase besitzt eine endliche Aufnahmekapazität für die Krümmungsrückgabe. Die Rückgaberate ist proportional zur verbleibenden Kapazität:
\begin{equation}
\frac{d\Omega_\Phi}{da} = k \cdot \left[1 - \left(\frac{\Omega_\Phi}{\Phi_0}\right)^{\!2}\right]
\label{eq:saturation_ode}
\end{equation}
Diese Gleichung beschreibt einen klassischen Sättigungsprozess der dynamischen Systemtheorie: Bei kleinem $\Omega_\Phi$ wächst das Potential nahezu linear (die ``Bremse'' löst sich mit voller Rate), bei $\Omega_\Phi \to \Phi_0$ tritt Sättigung ein (die maximale Kapazität ist erreicht, die Bremse vollständig gelöst). Die exakte Lösung von Gl.~\eqref{eq:saturation_ode} ist:
\begin{equation}
\Omega_\Phi(a) = \Phi_0 \cdot \tanh\!\big(k\cdot(a - a_{\mathrm{trans}})\big)
\end{equation}
wobei $a_{\mathrm{trans}}$ die Integrationskonstante (Übergangspunkt) darstellt. Der Sättigungsmechanismus ist in der Physik ubiquitär und tritt in formal identischer Form in zahlreichen Systemen auf:
\begin{itemize}
\item Ferromagnetismus: Spontane Magnetisierung $M(T) \sim \tanh(T_C/T)$
\item BCS-Supraleitung: Energielücke $\Delta(T) \sim \tanh(T_C/T)$
\item Solitonenphysik: Kink-Lösung $\phi(x) = \phi_0 \tanh(kx)$
\item Nichtlineare Optik: Sättigungsabsorption $\alpha(I) \propto 1/(1 + I/I_{\mathrm{sat}})$
\end{itemize}
Alle diese Systeme teilen die Eigenschaft eines geordneten Übergangs von einem Zustand in einen anderen mit endlicher Kapazität -- \textit{genau} das Verhalten, das der spieltheoretische Rahmen für das Nash-Gleichgewicht zwischen Nullraum und Raumzeitblase vorhersagt. Die $\tanh$-Form ist damit nicht postuliert, sondern aus dem zugrunde liegenden Mechanismus \textit{abgeleitet}. \textit{Mikroskopische Realisierung:} Die Sättigungs-ODE~\eqref{eq:saturation_ode} wird nicht lediglich auf phänomenologischer Ebene postuliert, sondern ergibt sich später mikroskopisch aus einem Pöschl-Teller-Potential in der effektiven Lagrange-Dichte (Paper~III \cite{Geiger2026c}): Das $\cosh^{-2}$-Potential erlaubt $\tanh$-artige Kink-Lösungen als universelle Mean-Field-Dynamik und liefert damit ein rigoroses feldtheoretisches Fundament für den Sättigungsmechanismus.

Zur Überprüfung der Robustheit wurden vier verschiedene Sättigungsfunktionen getestet (Abschnitt~\ref{subsec:funcforms}). Alle liefern $\Delta\chi^2 \approx -9$ bis $-12$ gegenüber $\Lambda$CDM -- die Daten ``sehen'' einen Sättigungsprozess, unabhängig von der exakten mathematischen Formulierung.

\subsection{Physikalische Interpretation der Parameter}
\label{subsec:interpretation}

\textbf{Frühe Zeiten} ($a \to 0$, $z \to \infty$): $\Omega_\Phi \to 0$. Die ``Bremse'' wirkt voll -- die Expansion folgt der Materiedominanz wie in $\Lambda$CDM. Es gibt keine dunkle Komponente.

\textbf{Übergangsepoche} ($a \approx a_{\mathrm{trans}}$, $z \approx 1{,}5$): $\Omega_\Phi$ steigt an. Die ``Bremse'' beginnt nachzulassen. Dies geschah vor etwa 10,3 Milliarden Jahren.

\textbf{Heute} ($a = 1$, $z = 0$): $\Omega_\Phi \to \Phi_0$. Der maximale Effekt ist erreicht; das Potential wirkt effektiv wie~$\Lambda$.

\subsection{Effektiver Zustandsgleichungsparameter}
\label{subsec:weff}

Der effektive Zustandsgleichungsparameter des Krümmungs-Rückgabepotentials ist:
\begin{equation}
w_{\mathrm{eff}}(a) = -1 - \frac{1}{3}\,\frac{d\ln\Omega_\Phi}{d\ln a}
\label{eq:weff}
\end{equation}

Seine Zeitentwicklung ist in Tabelle~\ref{tab:weff} dargestellt.

\begin{table}[H]
\centering
\caption{Zeitentwicklung des effektiven Zustandsgleichungsparameters $w_{\mathrm{eff}}(z)$ im Vergleich $\Lambda$CDM vs.\ CRM. Die $1\sigma$-Unsicherheiten stammen aus der MCMC-Posterioranalyse (Abschnitt~\ref{subsec:mcmc}).}
\label{tab:weff}
\begin{tabular}{ccccc}
\toprule
$z$ & $w$ ($\Lambda$CDM) & $w$ (CRM) & $1\sigma$-Bereich & $\Delta w$ \\
\midrule
0,0 & $-1{,}000$ & $-1{,}355$ & $[-1{,}371;\;-1{,}339]$ & $\mathbf{-0{,}355}$ \\
0,3 & $-1{,}000$ & $-1{,}433$ & $[-1{,}645;\;-1{,}355]$ & $\mathbf{-0{,}433}$ \\
0,5 & $-1{,}000$ & $-1{,}450$ & $[-1{,}730;\;-1{,}358]$ & $\mathbf{-0{,}450}$ \\
0,8 & $-1{,}000$ & $-1{,}456$ & $[-1{,}759;\;-1{,}359]$ & $\mathbf{-0{,}456}$ \\
1,0 & $-1{,}000$ & $-1{,}454$ & $[-1{,}749;\;-1{,}359]$ & $\mathbf{-0{,}454}$ \\
1,5 & $-1{,}000$ & $-1{,}444$ & $[-1{,}696;\;-1{,}357]$ & $\mathbf{-0{,}444}$ \\
2,0 & $-1{,}000$ & $-1{,}432$ & $[-1{,}644;\;-1{,}355]$ & $\mathbf{-0{,}432}$ \\
\bottomrule
\end{tabular}
\end{table}

Innerhalb des phänomenologischen Pantheon+-Tanh-Fits liegt die Effektivdiagnostik in Tabelle~\ref{tab:weff} konsistent unter $w=-1$. Dies sollte nicht als fundamentale Phantomfeld-Vorhersage gelesen werden: Es ist eine Diagnose des gefitteten Ansatzes und muss gegen die DESI-DR2-Präferenz für $w_0>-1$, $w_a<0$ sowie gegen die Lagrange-Formulierung von Paper~III geprüft werden. Der relevante Falsifikationstest ist daher, ob der Krümmungs-Relaxationssektor die späte Expansionsgeschichte und den späteren Beschleunigungsübergang ohne eine persistente physische Phantomkomponente reproduzieren kann.

% ===================================================================
% 4. NUMERISCHE TESTS
% ===================================================================
\section{Numerische Tests und Modellvergleich}
\label{sec:numerics}

\subsection{Flachheitsbedingung}
\label{subsec:flatness}

Um die Zahl freier Parameter zu reduzieren und physikalische Konsistenz zu gewährleisten, wird die Flachheitsbedingung
\begin{equation}
\Omega_m + \Omega_\Phi(a{=}1) = 1
\label{eq:flatness}
\end{equation}
auferlegt. Daraus folgt für die Amplitude:
\begin{equation}
\Phi_0 = \frac{(1 - \Omega_m)(1 + s)}{\tanh\!\big(k\cdot(1 - a_{\mathrm{trans}})\big) + s}
\end{equation}
Das CRM hat damit drei kosmologische Freiheitsgrade ($\Omega_m$, $k$, $a_{\mathrm{trans}}$) plus einen Nuisance-Parameter ($M$), also insgesamt vier effektive Parameter -- nur zwei mehr als $\Lambda$CDM.

\subsection{Datenbasis: Pantheon+}
\label{subsec:pantheonplus}

Der Test erfolgt gegen den Pantheon+-Datensatz \cite{Scolnic2022}, einen großen öffentlich verfügbaren Katalog spektroskopisch bestätigter Typ-Ia-Supernovae. Die Freigabe enthält 1701 Lichtkurven für 1550 eindeutige SNe~Ia; nach dem hier verwendeten $z > 0{,}01$-Schnitt zur Vermeidung von Pekuliargeschwindigkeits-Dominanz gehen 1590 SN-Ia-Lichtkurven-Datenpunkte in den Fit ein, im Rotverschiebungsbereich $z = 0{,}0102$ bis $z = 2{,}2614$. Als Observable dient die bias-korrigierte scheinbare B-Band-Helligkeit \texttt{m\_b\_corr}. Die Analyse wird sowohl mit diagonalen Fehlern als auch mit der vollen statistisch-systematischen Kovarianzmatrix (STAT+SYS) des Pantheon+-Datensatzes durchgeführt.

\subsection{Methodik}
\label{subsec:methodology}

\textbf{Distanzberechnung:} Die Leuchtkraftentfernung wird über eine kumulative Trapezregel auf einem feinen $z$-Gitter ($N = 2000$ Stützstellen) berechnet und auf die Daten-Rotverschiebungen interpoliert. Dieses Verfahren ist numerisch stabil (Fehler $< 10^{-5}$) und ermöglicht schnelle Evaluation während der Optimierung.

\textbf{Nuisance-Parameter:} Der absolute Helligkeitsoffset $M = M_B + 5\log_{10}(c/H_0) + 25$, der die absolute Helligkeit und die Hubble-Konstante absorbiert, wird analytisch marginalisiert:
\begin{equation}
M_{\mathrm{best}} = \frac{\sum_i w_i (m_i^{\mathrm{obs}} - \mu_i^{\mathrm{th}})}{\sum_i w_i}, \quad w_i = \sigma_i^{-2}
\end{equation}

\textbf{Optimierung:} Parameterbestimmung mittels \textit{Differential Evolution} (globaler evolutionärer Optimizer) mit anschließender L-BFGS-B-Verfeinerung (\textit{polish}).

\textbf{MCMC-Unsicherheiten:} Für das CRM (flach) werden Parameterunsicherheiten mittels \textit{emcee} \cite{ForemanMackey2013} bestimmt (32~Walkers, 3000~Schritte, 500~Burn-in). Die Posteriorverteilungen liefern $1\sigma$-Konfidenzintervalle für alle Parameter einschließlich der abgeleiteten Größen $\Phi_0$ und $z_{\mathrm{trans}}$.

\textbf{Modellselektion:} Neben $\chi^2$ werden das Akaike-Informationskriterium (AIC~$= \chi^2 + 2k$) und das Bayes-Informationskriterium (BIC~$= \chi^2 + k \ln n$) berechnet, wobei $k$ die Zahl effektiver Parameter und $n$ die Datenpunktanzahl ist. Zur Überprüfung auf Overfitting wird zusätzlich eine 5-Fold-Kreuzvalidierung durchgeführt. Der vollständige Analysecode ist öffentlich verfügbar.\footnote{\url{https://github.com/research-line/crm-cosmology}}

\subsection{Ergebnisse}
\label{subsec:results}

Es werden drei Modelle gefittet: flaches $\Lambda$CDM (2~Parameter), CRM mit Flachheitsbedingung (4~Parameter) und CRM ohne Einschränkung (5~Parameter).

\begin{table}[htbp]
\centering
\caption{Gefittete Parameter und Anpassungsgüte: $\Lambda$CDM vs.\ CRM gegen Pantheon+ (1590~SN-Ia-Datenpunkte nach dem $z>0{,}01$-Schnitt). Für das CRM~(flach) werden $1\sigma$-MCMC-Unsicherheiten angegeben.}
\label{tab:results}
\begin{tabular}{lccc}
\toprule
 & $\Lambda$CDM & CRM (flach) & CRM (frei) \\
\midrule
Freie Parameter $k$ & 2 & 4 & 5 \\
$\Omega_m$ & 0,244 & $0{,}368^{+0{,}025}_{-0{,}023}$ & 0,552 \\
$\Omega_\Lambda$ / $\Omega_\Phi(z{=}0)$ & 0,756 & 0,636 & 0,872 \\
$\Phi_0$ (abgeleitet) & -- & $0{,}988^{+0{,}615}_{-0{,}221}$ & 1,292 \\
$k$ (Übergangsschärfe) & -- & $1{,}44^{+1{,}22}_{-0{,}84}$ & 1,98 \\
$a_{\mathrm{trans}}$ ($z_{\mathrm{trans}}$) & -- & 0,75 (0,33) & 0,80 (0,25) \\
$\Omega_{\mathrm{total}}$ & 1,000 & 1,000 & 1,423 \\
\midrule
$\chi^2$ (diagonal) & 729,0 & 716,8 & 715,9 \\
$\chi^2$ (volle Kov.) & 1432,0 & 1420,8 & -- \\
$\chi^2/\mathrm{dof}$ & 0,459 & 0,452 & 0,452 \\
AIC (diagonal) & 733,0 & 724,8 & 725,9 \\
AIC (volle Kov.) & 1436,0 & 1428,8 & -- \\
BIC & 743,7 & 746,3 & 752,8 \\
\bottomrule
\end{tabular}
\end{table}

Das CRM mit Flachheitsbedingung zeigt $\Omega_m = 0{,}368 \pm 0{,}024$ (MCMC). \textit{Spannungshinweis:} Dieser Wert liegt um $\Delta\Omega_m = 0{,}053$ über dem Planck-CMB-Wert $\Omega_m^{\mathrm{Planck}} = 0{,}315 \pm 0{,}007$ \cite{Planck2020}; kombiniert mit der CRM-Unsicherheit entspricht das etwa $2{,}1\sigma$. Diese Abweichung ist kein kleiner Rundungseffekt. Sie kann bedeuten, dass das phänomenologische Potential $\Omega_\Phi(a)$ einen Teil der Materiedichte absorbiert, oder dass ein reiner SN-Ia-Fit $\Omega_m$ ohne CMB- und BAO-Daten nicht ausreichend bestimmt. Ein gemeinsamer Fit mit CMB und BAO ist nötig, um diese Spannung zu bewerten. Die gefittete Übergangsrotverschiebung $z_{\mathrm{trans}} = 0{,}33$ ($a_{\mathrm{trans}} = 0{,}75$) liegt bei späteren kosmischen Zeiten als theoretisch erwartet. Die Übergangsschärfe $k = 1{,}44^{+1{,}22}_{-0{,}84}$ beschreibt einen sanften Übergang mit breiter Posterior -- die Daten präferieren einen Übergang, lassen aber eine Bandbreite von Übergangsschärfen zu.

\textit{Hinweis zu $\chi^2/\mathrm{dof} < 1$:} Die reduzierten Chi-Quadrat-Werte ($\chi^2/\mathrm{dof} = 0{,}452$--$0{,}459$) liegen für beide Modelle deutlich unter~1{,}0. Dies ist charakteristisch für Pantheon+ SN-Ia-Analysen unter Verwendung der diagonalen Kovarianzmatrix, von der bekannt ist, dass sie individuelle Messunsicherheiten überschätzt (die statistische Streuung übersteigt die systematische Untergrenze). Der Effekt ist nicht modellspezifisch: Sowohl $\Lambda$CDM als auch CRM zeigen das gleiche $\chi^2/\mathrm{dof} \approx 0{,}46$, was bestätigt, dass die AIC/BIC-Differenzen nicht durch Fehlerüberschätzung getrieben werden, sondern echte Unterschiede in der Anpassungsgüte widerspiegeln. Die Analyse mit voller Kovarianzmatrix ($\chi^2/\mathrm{dof} \approx 0{,}91$, nahe bei eins) liefert ein besser kalibriertes Fehlermodell.

\textbf{Volle Kovarianzmatrix:} Die Wiederholung der Analyse mit der vollen statistisch-systematischen Kovarianzmatrix bestätigt die Ergebnisse: $\Delta\chi^2 = -11{,}2$ und $\Delta\mathrm{AIC} = -7{,}2$ (gegenüber $-12{,}2$ und $-8{,}2$ bei diagonalen Fehlern). Die leichte Reduktion erklärt sich durch die Berücksichtigung systematischer Korrelationen zwischen benachbarten Supernovae.

\subsection{Modellselektion}
\label{subsec:modelselection}

\begin{table*}[htbp]
\centering
\caption{Modellvergleich: CRM vs.\ $\Lambda$CDM. Negative Werte bevorzugen CRM.}
\label{tab:comparison}
\begin{tabular}{lcc}
\toprule
\textbf{Kriterium} & CRM (flach) vs.\ $\Lambda$CDM & CRM (frei) vs.\ $\Lambda$CDM \\
\midrule
$\Delta\chi^2$ & $\mathbf{-12{,}2}$ & $-13{,}1$ \\
$\Delta$AIC & $\mathbf{-8{,}2}$ & $-7{,}1$ \\
$\Delta$BIC & $+2{,}6$ & $+9{,}0$ \\
\midrule
5-Fold $\langle\chi^2/n\rangle$ & $\mathbf{0{,}4499}$ & $0{,}4498$ \\
$\Lambda$CDM: $\langle\chi^2/n\rangle$ & \multicolumn{2}{c}{$0{,}4519$} \\
\bottomrule
\end{tabular}
\end{table*}

\textbf{Interpretation:} Zwei Selektionskriterien bevorzugen das flache CRM direkt gegenüber $\Lambda$CDM, nämlich $\chi^2$ ($-12{,}2$) und AIC ($-8{,}2$); der rohe Kreuzvalidierungs-Mittelwert liegt für CRM leicht niedriger ($0{,}4499$ vs.\ $0{,}4519$). Das BIC, das zusätzliche Parameter strenger bestraft, bevorzugt dagegen $\Lambda$CDM ($\Delta\mathrm{BIC} = +2{,}6$). Nach der Kass-Raftery-Skala \cite{KassRaftery1995} liegt dies im Bereich positiver Evidenz gegen das komplexere Modell. Kreuzvalidierung ist der direkteste Overfitting-Test, aber ohne berichtete Fold-Standardabweichung darf die Differenz von 0{,}4\,\% nur als vergleichbare Generalisierung, nicht als entscheidende Evidenz für CRM gelesen werden.

\subsection{MCMC-Posterioranalyse}
\label{subsec:mcmc}

Die Parameterunsicherheiten werden mittels \textit{emcee} \cite{ForemanMackey2013} bestimmt (32~Walkers, 3000~Schritte, 500~Burn-in, Akzeptanzrate: 63\,\%). \textit{Konvergenz-Hinweis:} Formale Konvergenzdiagnostiken (Gelman--Rubin-$\hat{R}$-Statistik \cite{GelmanRubin1992}; effektive Stichprobengröße $N_\mathrm{eff}$) wurden in der ursprünglichen Analyse nicht berichtet. Die Akzeptanzrate von 63\,\% liegt im üblichen \textit{emcee}-Zielbereich (20--70\,\%), und die visuelle Inspektion der Walker-Trajektorien zeigte keine offensichtliche Nicht-Konvergenz; für eine publikationsreife Fassung ist aber eine quantitative Konvergenzzertifizierung (z.\,B.\ $\hat{R}<1{,}01$, $N_\mathrm{eff}/\mathrm{chain}>50$) empfohlen. Die Ergebnisse für das CRM~(flach) sind in Tabelle~\ref{tab:results} als $1\sigma$-Unsicherheiten angegeben. Die Posteriorverteilung von $\Omega_m$ ist nahezu gaußförmig mit $\Omega_m = 0{,}368^{+0{,}025}_{-0{,}023}$. Die Übergangsschärfe $k$ zeigt eine breite, asymmetrische Posterior ($k = 1{,}44^{+1{,}22}_{-0{,}84}$), was bedeutet, dass die Daten einen Übergang bevorzugen, aber die Schärfe weniger stark einschränken. Die abgeleiteten Größen $\Phi_0 = 0{,}988^{+0{,}615}_{-0{,}221}$ und $z_{\mathrm{trans}} = 0{,}35$ (MCMC-Posterior-Median; die Punktschätzung in Tabelle~\ref{tab:results} gibt $0{,}33$) stimmen innerhalb der breiten Posteriori von $a_{\mathrm{trans}}$ überein. Die aus dem MCMC berechneten $w(z)$-Konfidenzbänder (Tabelle~\ref{tab:weff}) zeigen, dass $w = -1$ für alle Rotverschiebungen außerhalb des $1\sigma$-Bereichs liegt.

\textbf{Interpretation von $\Omega_m = 0{,}368$:} Der CRM-Wert liegt über dem Planck-CMB-Wert ($\Omega_m^{\mathrm{Planck}} = 0{,}315 \pm 0{,}007$), was typisch für reine Supernova-Fits ist. Im CRM-Kontext hat diese Abweichung eine mögliche physikalische Deutung: Das Modell könnte einen Teil der in $\Lambda$CDM als ``Dunkle Energie'' klassifizierten Dichte als dynamischen Krümmungseffekt absorbieren. Da $\Omega_\Phi(a)$ bei frühen Zeiten gegen Null geht (im Gegensatz zu $\Omega_\Lambda = \mathrm{const.}$), muss $\Omega_m$ kompensatorisch höher ausfallen, um den Gesamt-Fit zu erhalten. Dies sollte jedoch nicht als unabhängige Stützung durch schwache Gravitationslinsen-Surveys gelesen werden: Lensing-Daten beschränken vor allem Kombinationen wie $S_8$ und besitzen eigene Spannungen mit Planck. Erst ein gemeinsamer SN+CMB+BAO+Lensing-Fit kann klären, ob der höhere SN-only-$\Omega_m$-Wert physikalisch zu interpretieren ist.

\subsection{Robustheit: Alternative Funktionalformen}
\label{subsec:funcforms}

Um zu überprüfen, ob die Ergebnisse von der spezifischen Wahl der $\tanh$-Form abhängen, werden vier verschiedene Sättigungsfunktionen unter identischen Bedingungen getestet:

\begin{table}[H]
\centering
\caption{Vergleich alternativer Funktionalformen für $\Omega_\Phi(a)$. Alle Modelle verwenden die Flachheitsbedingung und haben $k=4$ Parameter.}
\label{tab:funcforms}
\begin{tabular}{lcccc}
\toprule
\textbf{Funktionalform} & $\chi^2$ & AIC & $\Delta\chi^2$ vs.\ $\Lambda$CDM & $\Omega_m$ \\
\midrule
$\tanh$ (Standard-CRM) & 716,8 & 724,8 & $-12{,}2$ & 0,364 \\
Logistische Funktion & 717,5 & 725,5 & $-11{,}5$ & 0,368 \\
Error-Funktion (erf) & 716,7 & 724,7 & $-12{,}3$ & 0,367 \\
Potenzgesetz & 720,1 & 728,1 & $\phantom{0}-8{,}9$ & 0,364 \\
\bottomrule
\end{tabular}
\end{table}

\textbf{Ergebnis:} Alle vier Funktionalformen liefern $\Delta\chi^2 \approx -9$ bis $-12$ gegenüber $\Lambda$CDM. Die $\tanh$-Form ist weder die einzig mögliche noch die ``bestpassende'' -- sie repräsentiert eine robuste Klasse von Sättigungsfunktionen. Dies entkräftet den Einwand, die CRM-Ergebnisse seien von einer spezifischen Funktionswahl abhängig. Die $\Omega_m$-Werte konvergieren bei allen Formen auf $\approx 0{,}36$--$0{,}37$, was die physikalische Konsistenz unterstreicht.

\subsection{Phantom-Hintergrund und Stabilitätsvorbehalt}
\label{subsec:phantom}

Ein effektives $w<-1$ belegt für sich weder einen physikalischen Geist noch eine Singularität in endlicher Zeit. Für eine separat erhaltene Komponente mit konstantem $w<-1$ wächst die bekannte Skalierung $\rho\propto a^{-3(1+w)}$ mit $a$ und kann zu einem Big Rip führen \cite{Caldwell1998}. Der CRM-Hintergrundansatz unterscheidet sich hiervon wie folgt:

\begin{enumerate}
\item \textbf{Beschränkter Hintergrundbeitrag:} Die Sättigungsgleichung (Gl.~\ref{eq:saturation_ode}) liefert $\Omega_\Phi(a)\to\Phi_0$. Im verwendeten Friedmann-Ansatz bleibt dieser Beitrag endlich.
\item \textbf{Asymptotischer de-Sitter-Hintergrund:} Bei abklingendem Materiebeitrag nähert sich $H$ einer endlichen Konstante und $w_{\mathrm{eff}}\to-1$. Dies ist eine Aussage über den homogenen Hintergrund, kein Beweis seiner Störungsstabilität.
\item \textbf{Kein Big Rip in diesem Ansatz:} Ein spätzeitlich beschränktes $H$ schließt eine Divergenz von $a$ in endlicher Eigenzeit aus. Diese Aussage gilt unter den Annahmen des Hintergrundansatzes und ist enger als die Abwesenheit sämtlicher Phantom-Pathologien.
\item \textbf{Effektive geometrische Beschreibung:} Eine geometrischen Termen zugeordnete effektive Zustandsgleichung ist nicht mit dem Energie-Impuls-Tensor eines physikalischen Phantomfelds identisch \cite{Sotiriou2010}. Diese Unterscheidung ersetzt keine Prüfung der kinetischen Vorzeichen, Gradienteninstabilitäten und CRM-Feldgleichungen.
\end{enumerate}

\begin{quote}
\textit{Der Sättigungsansatz schließt unter seinen Voraussetzungen einen Hintergrund-Big-Rip aus. Geisterfreiheit oder Störungsstabilität folgen daraus allein nicht.}
\end{quote}

\subsection{Dezelerationsparameter und $H_0$-Implikationen}
\label{subsec:deceleration}

\textbf{Dezelerationsparameter $q(z)$:} Der Dezelerationsparameter liefert eine zusätzliche, unabhängig testbare Vorhersage. Im CRM tritt der Übergang von gebremster zu beschleunigter Expansion ($q = 0$) bei $z_{\mathrm{acc}} = 0{,}52$ auf -- deutlich später in kosmischer Zeit als im $\Lambda$CDM-Modell ($z_{\mathrm{acc}} = 0{,}84$). Zudem sagt das CRM eine stärkere heutige Beschleunigung voraus: $q_0^{\mathrm{CRM}} = -0{,}81$ gegenüber $q_0^{\Lambda\mathrm{CDM}} = -0{,}63$.

\textbf{Implikation für die Strukturbildung:} Ein späterer Beginn der kosmischen Beschleunigung ($z_{\mathrm{acc}} = 0{,}52$ statt $0{,}84$) bedeutet, dass gravitativ gebundene Strukturen (Galaxienhaufen, großräumige Filamente) \textit{länger ungestört wachsen konnten}, bevor die Beschleunigung das Wachstum unterdrückte. Die materiedominierte Ära -- in der die Gravitation das Strukturwachstum antreibt -- dauerte im CRM signifikant länger an als im Standardmodell.

\textbf{Empirische Evidenz für frühe massive Strukturen:} Mehrere unabhängige Beobachtungen setzen das $\Lambda$CDM-Modell in Bezug auf die Strukturbildung unter erheblichen Druck:
\begin{enumerate}
\item \textbf{JWST ``Universe Breakers'':} Das James Webb Space Telescope hat Galaxien bei $z > 7$ (ca.\ 500--700\,Myr nach dem Urknall) entdeckt, die weitaus massereicher sind als von hierarchischer Strukturbildung im $\Lambda$CDM erlaubt \cite{Labbe2023}. Boylan-Kolchin \cite{BoylanKolchin2023} zeigt quantitativ, dass die stellare Massendichte dieser Objekte die verfügbare Baryonen-Budget innerhalb von $\Lambda$CDM-Halos übersteigt -- ein ``Timing-Problem'' der Strukturbildung.
\item \textbf{El~Gordo (ACT-CL~J0102$-$4915):} Dieser extrem massereiche Galaxienhaufen bei $z \approx 0{,}87$ mit Masse $M_{200} \approx 2{,}1 \times 10^{15}\,M_\odot$ stellt eine $> 6\sigma$-Spannung mit $\Lambda$CDM dar, da ein Objekt dieser Masse bei diesem Alter mit der beobachteten Kollisionsgeschwindigkeit statistisch nahezu unmöglich ist \cite{Asencio2023}.
\item \textbf{Protocluster SPT2349$-$56:} Bereits 1,4\,Mrd.\ Jahre nach dem Urknall ($z = 4{,}3$) zeigt dieser Protocluster mindestens 14 gasreiche Galaxien mit einer Gesamtsternentstehungsrate von $\sim\!6500\,M_\odot$/yr -- weit reifer als von $\Lambda$CDM vorhergesagt \cite{Miller2018}.
\end{enumerate}

Das CRM bietet für alle drei Beobachtungen eine natürliche Erklärung: Da die Beschleunigung erst bei $z_{\mathrm{acc}} = 0{,}52$ einsetzt (statt $z_{\mathrm{acc}} = 0{,}84$), stand gravitativ gebundenen Strukturen mehr kosmische Zeit für ungestörtes Wachstum zur Verfügung. Die CRM-Vorhersage ist direkt testbar durch Galaxienhaufen-Zählungen und schwache Gravitationslinsen-Surveys (Euclid, Vera C.\ Rubin Observatory).

\begin{table}[H]
\centering
\caption{Dezelerationsparameter $q(z)$: $\Lambda$CDM vs.\ CRM.}
\label{tab:qz}
\begin{tabular}{cccc}
\toprule
$z$ & $q$ ($\Lambda$CDM) & $q$ (CRM) & $\Delta q$ \\
\midrule
0,0 & $-0{,}634$ & $-0{,}805$ & $-0{,}171$ \\
0,5 & $-0{,}217$ & $-0{,}017$ & $+0{,}200$ \\
1,0 & $+0{,}082$ & $+0{,}321$ & $+0{,}240$ \\
2,0 & $+0{,}346$ & $+0{,}467$ & $+0{,}121$ \\
\bottomrule
\end{tabular}
\end{table}

\textbf{$H_0$-Implikationen:} Der Nuisance-Parameter $M$ absorbiert sowohl die absolute Helligkeit $M_B$ als auch $H_0$ über die Beziehung $M = M_B + 5\log_{10}(c/H_0) + 25$. Unter Verwendung der SH0ES-Eichung ($M_B = -19{,}253$) liefert das CRM $H_0 = 76{,}1\,\mathrm{km/s/Mpc}$ gegenüber $H_0 = 75{,}5\,\mathrm{km/s/Mpc}$ im $\Lambda$CDM -- ein Unterschied von $\Delta H_0 = +0{,}5\,\mathrm{km/s/Mpc}$. Die $H_0$-Spannung wird durch das CRM allein nicht gelöst, da der Unterschied innerhalb der Messunsicherheit liegt. Eine direkte Auflösung erfordert die Kombination mit CMB- und BAO-Daten.

% ===================================================================
% 5. VERGLEICH MIT ALTERNATIVEN
% ===================================================================
\section{Vergleich mit alternativen Modellen}
\label{sec:alternatives}

\subsection{$\Lambda$CDM (Standardmodell)}

Das $\Lambda$CDM-Modell ist extrem einfach ($w = -1$, konstant, zwei kosmologische Parameter) und passt alle aktuellen Daten gut. Es leidet jedoch unter dem Kosmologische-Konstante-Problem und dem Koinzidenz-Problem \cite{Weinberg1989}.

\subsection{Quintessenz}

Quintessenz-Modelle \cite{Caldwell1998} postulieren ein dynamisches Skalarfeld~$\phi$ mit zeitabhängigem Zustandsgleichungsparameter. Sie können das Koinzidenz-Problem mildern, erfordern aber ein neues Feld und dessen Potential~$V(\phi)$ mit vielen freien Parametern.

\subsection{Modifizierte Gravitation: $f(R)$-Theorien}

$f(R)$-Gravitationstheorien \cite{Starobinsky1980, Sotiriou2010} ersetzen den Ricci-Skalar~$R$ in der Einstein-Hilbert-Wirkung durch eine allgemeinere Funktion. Sie bieten eine geometrische Erklärung ohne Dunkle Energie, sind jedoch mathematisch komplex und zum Teil inkonsistent mit Beobachtungen (Gravitationslinsen, CMB). Das am weitesten untersuchte spätzeitliche $f(R)$-Modell, Hu-Sawicki \cite{HuSawicki2007}, erfüllt die Constraints des Sonnensystems durch den Chamäleon-Mechanismus und erzeugt gleichzeitig eine viable kosmische Beschleunigung; es dient als wichtiger Vergleichsmaßstab für die Lagrange-Analyse in Paper~III \cite{Geiger2026c}. Der rein quadratische $R^2$-Term des CRM teilt diese Abschirmungseigenschaft \textit{nicht}, wie dort ausgeführt.

\subsection{Emergente Gravitation nach Verlinde}

Verlinde \cite{Verlinde2011, Verlinde2017} schlägt vor, dass die Gravitation keine fundamentale Kraft, sondern ein emergentes, entropisches Phänomen ist. In de-Sitter-Räumen führt die mit dem kosmologischen Horizont assoziierte Entropie zu einer zusätzlichen ``dunklen'' Gravitationskraft, die das Verhalten von Galaxien ohne Dunkle Materie erklären könnte. In einem eng verwandten Ansatz leitet Padmanabhan \cite{Padmanabhan2012} die kosmische Expansion als emergentes Phänomen aus der Differenz zwischen Oberflächen- und Bulk-Freiheitsgraden ab, wobei die kosmologische Konstante aus dem Bestreben der Raumzeit entsteht, holographische Äquipartition zu erreichen -- eine Perspektive, die strukturelle Gemeinsamkeiten mit dem Gleichgewichtsrahmenwerk des CRM aufweist.

\subsection{Finsler-Gravitation}

Pfeifer et al.\ \cite{Pfeifer2025} erweitern die Allgemeine Relativitätstheorie durch Finsler-Geometrie, in der die Metrik nicht nur von der Position, sondern auch von der Geschwindigkeit abhängt:
\begin{equation}
g_{\mu\nu}(x, y) = \frac{1}{2}\,\frac{\partial^2 F^2}{\partial y^\mu \partial y^\nu}, \quad y = \frac{dx}{d\lambda}
\end{equation}
Die resultierende Finsler-Friedmann-Gleichung erzeugt selbst im Vakuum eine exponentielle Expansion -- ohne kosmologische Konstante.

\subsection{Cosmological Teleodynamics}

Trivedi und Venkatasubramanian \cite{Trivedi2025} formulieren eine spieltheoretische Kosmologie, die erstaunliche Parallelen zum hier vorgestellten Ansatz aufweist. Ihre \textit{Cosmological Teleodynamics} beschreibt das Universum als ``riesiges Potentialspiel'', das sich einem kontinuierlichen Nash-Gleichgewicht annähert. Die kosmische Beschleunigung erscheint als ``statistisch emergenter Effekt dynamischen Gedächtnisses in einem selbstgravitierenden Medium'' -- eine Formulierung, die konzeptuell dem ``geometrischen Gedächtnis'' des CRM entspricht.

\subsection{Modelle mit dynamischer Raumkrümmung}
\label{subsec:dynamische_kruemmung}

Eine weitere Klasse von Ansätzen ohne Dunkle Energie führt die spätzeitliche kosmische Beschleunigung auf eine zeitabhängige effektive Raumkrümmung statt auf eine kosmologische Konstante zurück. Stichel \cite{Stichel2016} formuliert ein relativistisches ``Darkon''-Fluid-Modell auf einer konform flachen Raumzeitgeometrie, dessen dynamische Raumkrümmung einen Übergang von Abbremsung zu spätzeitlicher Beschleunigung ohne explizite Dunkle-Energie-Komponente antreibt. Jüngst untersuchen Esteban-Gutiérrez et al.\ \cite{EstebanGutierrez2025} einen FLRW-Hintergrund, in dem der Raumkrümmungsparameter einen glatten phänomenologischen Übergang zwischen zwei asymptotischen Krümmungsregimen durchläuft und Cosmic-Chronometer- sowie Typ-Ia-Supernova-Daten ohne $\Lambda$ beschreibt. Obwohl diese Modelle mit dem CRM das übergeordnete Ziel teilen, die Spätzeitbeschleunigung geometrisch ohne Vakuumenergiekonstante zu erklären, ist ihr physikalischer Mechanismus strukturell verschieden: Modelle mit dynamischer Raumkrümmung modifizieren den Friedmann-Raumkrümmungsterm~$k(t)/a^2$, während das CRM einen räumlich flachen FLRW-Hintergrund ($k = 0$) beibehält und die effektive Sättigungsdynamik $\Omega_\Phi(a) = \Phi_0 \tanh(k a)$ aus Krümmungsgedächtnis-Relaxation und thermodynamischer Äquipartition ableitet.

\subsection{Synoptischer Vergleich}
\label{subsec:synopse}

\begin{table*}[tbp]
\centering
\caption{Synoptischer Vergleich kosmologischer Modelle ohne Dunkle Energie.}
\label{tab:synopse}
\begin{tabular}{llll}
\toprule
\textbf{Eigenschaft} & \textbf{CRM} & \textbf{Finsler} & \textbf{Teleodynamics} \\
\midrule
Theor.\ Basis & Standard-ART + Potential & Finsler-Geometrie & Stat.\ Mechanik + Spieltheorie \\
Mechanismus & Nachlassende ``Bremse'' & Geschwindigkeitsabh.\ Metrik & Dynamisches Gedächtnis \\
Dunkle Energie & Nicht nötig & Nicht nötig & Nicht nötig \\
Empirischer Test & Pantheon+ (1590 Datenpkt., $\Delta\chi^2{=}{-}12$) & Noch ausstehend & Qualitativ \\
Vorhersage & $w(z)$ Zeitvariation & Exp.\ Expansion & Nash-Konvergenz \\
Komplexität & Gering (4 Param.) & Hoch & Mittel \\
\bottomrule
\end{tabular}
\end{table*}

% ===================================================================
% 6. KOMPLEMENTARITÄT UND VEREINIGUNG
% ===================================================================
\section{Komplementarität und mögliche Vereinigung}
\label{sec:complementarity}

\subsection{Drei Modelle, eine Einsicht}

Alle drei Ansätze -- CRM, Finsler-Gravitation und Cosmological Teleodynamics -- teilen eine fundamentale Einsicht:
\begin{quote}
\textit{``Die beschleunigte Expansion ist kein neues `Ding', sondern eine Eigenschaft der Geometrie bzw.\ der statistischen Struktur des Universums selbst.''}
\end{quote}

\subsection{Hypothese: CRM als effektive Beschreibung}

Eine faszinierende Möglichkeit besteht darin, dass die drei Modelle verschiedene Aspekte desselben Phänomens beschreiben. In Analogie zur Beziehung zwischen Thermodynamik und Statistischer Mechanik könnte gelten:

\begin{itemize}
\item \textbf{Finsler-Gravitation} (mikroskopisch, fundamental): Alle Momente der 1-Partikel-Verteilungsfunktion tragen zur Gravitation bei.
\item \textbf{CRM} (makroskopisch, phänomenologisch): Das zeitabhängige Potential $\Phi(a)$ kodiert effektiv den Beitrag der höheren Momente.
\item \textbf{Teleodynamics} (systemisch, spieltheoretisch): Die Nash-Gleichgewichtsdynamik beschreibt die globale Optimierung.
\end{itemize}

Mathematisch lässt sich diese Komplementarität als hierarchische Beziehung darstellen:
\begin{equation}
\begin{split}
\underbrace{G_{\mu\nu}^{\mathrm{Finsler}}(x, y)}_{\substack{\text{Finsler-Gravitation}\\\text{(mikroskopisch)}}}
&\;\xrightarrow{\;\langle\cdot\rangle_{\mathrm{eff}}\;}
\underbrace{G_{\mu\nu} + 8\pi G\,\Omega_\Phi(a)\,g_{\mu\nu}}_{\substack{\text{CRM: modifizierte}\\\text{Einstein-Gleichung}}} \\
&\;\xleftarrow{\;\delta\Phi/\delta s_i = 0\;}
\underbrace{\max_{s_M, s_T} \Phi(s_M, s_T)}_{\substack{\text{Teleodynamics:}\\\text{Nash-Gleichgewicht}}}
\end{split}
\label{eq:complementarity}
\end{equation}
wobei der linke Pfeil die effektive Mittelung über die geschwindigkeitsabhängigen Freiheitsgrade der Finsler-Geometrie beschreibt und der rechte Pfeil die Variationsbedingung des spieltheoretischen Potentials darstellt, die das zeitliche Verhalten von $\Omega_\Phi(a)$ festlegt.

% ===================================================================
% 7. TESTBARKEIT UND VORHERSAGEN
% ===================================================================
\section{Testbarkeit und Vorhersagen}
\label{sec:testability}

\subsection{Beobachtbare Signaturen}

\textbf{1.~Zustandsgleichungs-Diagnostik:} Der phänomenologische Pantheon+-Fit liefert $|\Delta w| \approx 0{,}4$ relativ zu $\Lambda$CDM, aber DESI DR2 verschiebt das physische Ziel weg von persistenter Phantomhaftigkeit hin zu einer $w_0>-1$, $w_a<0$-Trajektorie. Die ESA-Mission Euclid \cite{Euclid2024}, das Nancy Grace Roman Space Telescope und DESI können prüfen, ob diese Entwicklung besser durch Krümmungsrelaxation oder durch eine generische $w_0w_a$CDM-Parametrisierung beschrieben wird.

\textbf{2.~Dezelerationsparameter:} Das CRM sagt einen früheren Übergang zur beschleunigten Expansion voraus ($z_{\mathrm{acc}} = 0{,}52$ vs.\ $\Lambda$CDM: $z_{\mathrm{acc}} = 0{,}84$) sowie eine stärkere heutige Beschleunigung ($q_0 = -0{,}81$ vs.\ $-0{,}63$). Dies ist unabhängig von $w(z)$ testbar.

\textbf{3.~Strukturwachstum:} Eine modifizierte Wachstumsrate $f \cdot \sigma_8$ ist vorhergesagt, messbar durch schwache Gravitationslinsen und Galaxienhaufen-Zählungen. Erste empirische Hinweise liefern bereits die JWST-``Universe Breakers'' bei $z > 7$ \cite{Labbe2023}, die El-Gordo-Anomalie ($> 6\sigma$ Spannung mit $\Lambda$CDM; \cite{Asencio2023}) und unerwartet reife Protocluster bei $z > 4$ \cite{Miller2018} -- allesamt konsistent mit der CRM-Vorhersage einer verlängerten Wachstumsphase.

\textbf{4.~CMB-Integraleffekte:} Ein modifizierter ISW-Effekt (\textit{Integrated Sachs-Wolfe}) in CMB-Temperatur-Kreuzkorrelationen.

\subsection{Zukünftige Missionen}

\begin{table}[H]
\centering
\caption{Relevante Beobachtungsmissionen für den CRM-Test.}
\label{tab:missions}
\small
\begin{tabular}{llcl}
\toprule
\textbf{Mission} & \textbf{Start} & $\sigma(w)$ & \textbf{Relevanz} \\
\midrule
Euclid & 2023 & ${\approx}\,0{,}02$ & BAO + Linsen \\
Roman & 2026--2027 & ${\approx}\,0{,}03$ & SN bis $z \approx 2$ \\
DESI & 2021-- & ${\approx}\,0{,}04$ & BAO + Wachstum \\
\bottomrule
\end{tabular}
\end{table}

\subsection{Unterscheidbarkeit der Modelle}

\begin{table}[H]
\centering
\caption{Vergleich der Vorhersagen: $\Lambda$CDM vs.\ CRM (Pantheon+-Fit). Die $1\sigma$-Unsicherheiten stammen aus der MCMC-Analyse.}
\label{tab:predictions}
\begin{tabular}{lcc}
\toprule
\textbf{Eigenschaft} & $\Lambda$CDM & CRM \\
\midrule
$w(z{=}0)$ & $-1{,}000$ & $-1{,}36 \pm 0{,}02$ \\
$w(z{=}0{,}5)$ & $-1{,}000$ & $-1{,}45^{+0{,}09}_{-0{,}28}$ \\
$w(z{=}1)$ & $-1{,}000$ & $-1{,}45^{+0{,}10}_{-0{,}30}$ \\
$w(z{=}2)$ & $-1{,}000$ & $-1{,}43^{+0{,}08}_{-0{,}21}$ \\
Zeitvariation & Keine & Ja (durchgehend $w < -1$) \\
$\Delta w$ (messbar) & -- & $\approx -0{,}4$ \\
$q_0$ (heute) & $-0{,}63$ & $-0{,}81$ \\
$z_{\mathrm{acc}}$ (Übergang) & 0,84 & 0,52 \\
\bottomrule
\end{tabular}
\end{table}

% ===================================================================
% 8. DISKUSSION
% ===================================================================
\section{Diskussion}
\label{sec:discussion}

\subsection{Stärken des Ansatzes}

\begin{enumerate}
\item \textbf{Konzeptuelle Eleganz:} Keine neue Energieform erforderlich; die Beschleunigung ist eine ``nachlassende Einschränkung'', kein ``neuer Antrieb''.
\item \textbf{Spieltheoretische Fundierung:} Die Emergenz physikalischer Gesetze aus Gleichgewichtsbedingungen bietet einen neuartigen Erklärungsrahmen, der durch die unabhängige Arbeit von Trivedi und Venkatasubramanian \cite{Trivedi2025} gestützt wird.
\item \textbf{Empirische Validierung:} Das CRM passt 1590 Pantheon+ SN-Ia-Datenpunkte nach $\chi^2$ und AIC besser als $\Lambda$CDM ($\Delta\chi^2 = -12{,}2$, $\Delta\mathrm{AIC} = -8{,}2$); die Kreuzvalidierung ist mit $0{,}4499$ vs.\ $0{,}4519$ vergleichbar, aber ohne Fold-Varianz nicht entscheidend.
\item \textbf{Testbarkeit:} Spezifische, quantitative Vorhersagen für $w(z)$ und $z_{\mathrm{acc}}$, die innerhalb einer Dekade überprüfbar sind.
\item \textbf{Empirische Unterstützung:} Die CRM-Vorhersage einer verlängerten Wachstumsphase ($z_{\mathrm{acc}} = 0{,}52$) bietet eine natürliche Erklärung für die JWST-``Early Galaxy Tension'' \cite{Labbe2023, BoylanKolchin2023}, die El-Gordo-Anomalie \cite{Asencio2023} und unerwartet reife Protocluster \cite{Miller2018}.
\item \textbf{Konvergenz unabhängiger Ansätze:} CRM, Finsler-Gravitation und Cosmological Teleodynamics kommen unabhängig zum selben Schluss: Dunkle Energie ist nicht notwendig.
\item \textbf{Reproduzierbarkeit:} Analysecode und Daten sind öffentlich verfügbar (\url{https://github.com/research-line/crm-cosmology}).
\end{enumerate}

\subsection{Limitationen und offene Fragen}

\begin{enumerate}
\item \textbf{Phänomenologischer Charakter:} Das CRM ist keine fundamentale Theorie. Obwohl die $\tanh$-Form als exakte Lösung der Sättigungs-ODE~\eqref{eq:saturation_ode} motiviert werden kann und vier alternative Funktionalformen vergleichbare Ergebnisse liefern (Abschnitt~\ref{subsec:funcforms}), steht eine Herleitung aus einer fundamentalen Quantengleichung noch aus.
\item \textbf{Parameterfreiheit und BIC-Strafe:} Vier effektive Parameter gegenüber zwei in $\Lambda$CDM führen zu einem BIC-Nachteil ($\Delta\mathrm{BIC} = +2{,}6$), der nach Kass--Raftery positive Evidenz gegen das komplexere Modell bedeutet. Die Kreuzvalidierung relativiert den Overfitting-Verdacht, beweist aber ohne Fold-Varianzen keinen Vorteil. Eine ehrliche Kurzfassung lautet: AIC bevorzugt CRM, BIC bevorzugt $\Lambda$CDM, Kreuzvalidierung ist inkonklusiv; entscheidend werden direkte CPL-, CMB- und BAO-Vergleiche.
\item \textbf{Phantom-Bereich:} Das gefittete effektive $w<-1$ und der beschränkte Hintergrund werden in Abschnitt~\ref{subsec:phantom} diskutiert. Das Ausbleiben eines Big Rip in diesem Ansatz beweist weder Geisterfreiheit noch Störungsstabilität; beides bleibt gesondert an der tatsächlichen CRM-Wirkung zu prüfen \cite{Sotiriou2010}.
\item \textbf{Modellvergleich:} Der Vergleich in dieser Arbeit beschränkt sich auf $\Lambda$CDM. Ein fairerer Referenzpunkt für das Phantom-Kreuzungsverhalten wäre die $w_0w_a$CDM-Parametrisierung (CPL) \cite{Chevallier2001, Linder2003}, die ebenfalls $w < -1$ erlaubt. Ein solcher Vergleich wird in Paper~II präsentiert.
\item \textbf{Offene Tests:} CMB-Vorhersagen (ISW-Effekt, CMB-Leistungsspektrum), BAO-Signaturen und Gravitationslinsen-Effekte müssen noch berechnet werden. Die Analyse mit der vollen Kovarianzmatrix (Abschnitt~\ref{subsec:results}) bestätigt jedoch die Ergebnisse der diagonalen Analyse.
\item \textbf{Mikroskopische Basis:} Was ist $\Phi$ auf Quantenebene? Die Verbindung zu einer Theorie der Quantengravitation steht aus. Die mögliche Beziehung zur Finsler-Gravitation (Abschnitt~\ref{sec:complementarity}) könnte hier eine Brücke schlagen.
\item \textbf{$H_0$-Spannung:} Die $H_0$-Analyse (Abschnitt~\ref{subsec:deceleration}) zeigt $\Delta H_0 = +0{,}5\,$km/s/Mpc zwischen CRM und $\Lambda$CDM -- zu gering, um die $H_0$-Spannung zu lösen. Eine Auflösung erfordert die Kombination mit CMB- und BAO-Daten.
\item \textbf{DESI-DR2-Spannung und Lagrange-Theorie:} Der phänomenologische Pantheon+-Fit dieser Arbeit liefert $w_{\mathrm{eff}}< -1$ als phantomartige Diagnostik. DESI DR2 \cite{DESI2025} bevorzugt im $w_0w_a$CDM-Modell jedoch Lösungen im Quadranten $w_0>-1$, $w_a<0$; die veröffentlichte DESI+CMB+Union3-Zeile gibt $w_0=-0{,}667\pm0{,}088$ und $w_a=-1{,}09^{+0{,}31}_{-0{,}27}$ an, während Pantheon+ und DESY5 andere, aber ebenfalls nicht-phantomhafte $w_0$-Werte liefern. Diese Spannung betrifft die ursprüngliche Tanh-Lesart direkt und wird quantitativ an die Lagrange-Formulierung in Paper~III~\cite{Geiger2026c} delegiert: Der $f(R)=R+2\gamma R^2$-Sektor soll $w=-1$ von oben annähern, statt eine persistente physische Phantomkomponente zu verlangen. Der Wert $w=-1{,}355$ ist daher eine gefittete Effektivdiagnostik, keine Vorhersage der fundamentalen Theorie. Rekonstruierte Skalar-Tensor-Theorien \cite{Efstratiou2026} liefern ein modellabhängiges Beispiel, in dem nichtminimale Kopplungen $F(\phi)R$ eine effektive Phantomkreuzung ermöglichen und die abgeleitete $H_0$-Spannung verringern. Dies begründet weder ein allgemeines No-go-Theorem für Phantomkreuzungen in metrischer $f(R)$-Gravitation noch eine stabile Kreuzung im CRM. Hintergrundentwicklung und Störungsstabilität müssen für die tatsächliche CRM-Wirkung geprüft werden. Kovacs~\cite{Kovacs2026} ist in dieselbe Richtung relevant, weil die vorgeschlagene kovariante Erweiterung $w\geq -1$ als Kandidatenfolge der Curvature-Occupation-Kinetik aufführt; das Manuskript kennzeichnet die kosmologische Erweiterung aber ausdrücklich als konditional und noch nicht als abgeschlossene Störungstheorie. $w\geq -1$ ist daher das natürliche CRM-Testziel, kein bereits durch Paper~I bewiesenes Resultat.
\end{enumerate}

\subsection{Philosophische Implikationen}

Falls das CRM (oder ein verwandtes Modell) bestätigt wird, hätte dies tiefgreifende Konsequenzen:

\begin{itemize}
\item \textbf{Dunkle Energie ist kein ``Ding'':} Sie wäre eine geometrische Erinnerung, kein physisches Feld.
\item \textbf{Das Universum ``weiß'' von seinem Anfang:} Die Geometrie besitzt ein ``Gedächtnis''.
\item \textbf{Paradigmenwechsel:} Von ``Was treibt die Beschleunigung an?'' zu ``Warum bremste die Expansion früher?''
\end{itemize}

Dies wäre vergleichbar mit dem Übergang von ``Was treibt die Planeten an?'' (Ptolemäus: Sphären) zu ``Wie bewegen sich Planeten in der Geometrie des Raumes?'' (Kepler, Newton, Einstein).

\subsection{Hierarchische Nash-Gleichgewichte}
\label{subsec:hierarchical_nash}

Das Zweispielermodell aus Abschnitt~\ref{sec:gametheory} (Nullraum vs.\ Raumzeitblase) stellt nur die erste Ebene einer emergenten Hierarchie dar. Paper~IV \cite{Geiger2026d} zeigt, dass das Nash-Gleichgewicht des kosmologischen Spiels \emph{Tochterdynamik} erzeugt: Der skalare Sektor (Scalaron), der den kosmologischen Hintergrund liefert, reicht allein nicht aus, um dissipative Strukturen (Galaxien) auf galaktischer Skala zu stabilisieren. Eine zweite Gleichgewichtsebene -- ein zeitartiges Einheitsvektorfeld $A_\mu$, abgeleitet aus dem Prinzip maximaler Entropieproduktion -- wird benötigt, um flache Rotationskurven und die Radiale Beschleunigungsrelation zu erzeugen. Die hierarchische Struktur lautet somit:
\begin{quote}
\textit{Nullraum} $\to$ \textit{Skalarer Sektor} (Paper~III) $\to$ \textit{Vektorsektor} (Paper~IV) $\to$ \textit{Standardphysik}.
\end{quote}
Jede Ebene ist ein Nash-Gleichgewicht, das bei Erschöpfung den nächsten Freiheitsgrad erfordert. Der Zweispieler-Rahmen dieses Papers liefert die konzeptionelle Vorlage; Papers~III und~IV stellen die Lagrange-Realisierung bereit.

\subsection{Deduktive Struktur des Modells}
\label{subsec:deductive}

Ein Unterscheidungsmerkmal des CRM ist, dass sein Lagrangian nicht \textit{ad hoc} aus unabhängigen phänomenologischen Termen zusammengesetzt, sondern aus einer deduktiven Kette abgeleitet ist, die von zwei Annahmen ausgeht und bei testbaren Vorhersagen endet:

\begin{enumerate}
\item \textbf{Axiom~1 (Anfangsbedingung):} Ein metastabiles Quantenvakuum (Nullraum) lässt eine große Fluktuation zu, die einen Konzentrationsgradienten erzeugt -- die entstehende Raumzeit.
\item \textbf{Axiom~2 (Entwicklungsprinzip):} Das System entwickelt sich entlang des Pfades, der die Entropieproduktion maximiert, während die strukturelle Integrität des Nullraums erhalten bleibt (äquivalent: Prinzip der kleinsten Wirkung unter der Randbedingung der Muttersystem-Stabilität; siehe Abschnitt~\ref{subsec:thermodynamic}).
\end{enumerate}

Aus diesen beiden Annahmen führt folgende Deduktionskette -- hier rückblickend dargestellt; ihre strenge Herleitung ist Gegenstand von Paper~III -- zur CRM-Wirkung (Paper~III, Gl.~(8)):
\begin{itemize}
\item Axiom~2 erfordert eine \textit{kontrollierte, zeitlich gestreckte} Neutralisation des Gradienten $\to$ Einstein-Hilbert-Term $R/(16\pi G)$ (konservative Feldgleichungen als notwendige Bedingung).
\item Der Gradient hat endliche Größe, die Rückkehr darf das Vakuum nicht destabilisieren $\to$ \textit{Sättigung} der Rückgaberate $\to$ Pöschl-Teller-Skalarfeld mit $\tanh$-Dynamik.
\item Minimalität (Geistfreiheit + Renormierbarkeit auf Ein-Schleifen-Niveau) selektiert eindeutig $R^2$ als führende Gravitationskorrektur $\to$ Starobinsky-Term $\gamma R^2$.
\item Das Skalaron ($R^2$-Freiheitsgrad) erzeugt einen $a^{-\beta}$-Potenzgesetzterm, der die Rolle der Teilchen-Dunklen-Materie auf kosmologischer Ebene übernimmt; die Spurkopplung an $T_{\mu\nu}$ unterdrückt diesen Term automatisch während der Strahlungsära (BBN-Schutz).
\item Auf galaktischen Skalen verstärkt dasselbe $R^2$-Skalaron die Gravitation um $\mu \to 4/3$ auf Sub-Compton-Skalen (Paper~III), koinzidierend mit dem MOND-Phasenraumfaktor $V_3/V_2 = 4/3$ (Paper~II).
\item Die Sonnensystem-Kompatibilität folgt dagegen \textit{nicht} aus dieser Kette. Für ein rein quadratisches $f(R)$ ist $f_{RR}$ konstant, die Skalaron-Masse also krümmungs- und folglich dichteunabhängig, und es wirkt keine Chamäleon-Abschirmung. Eine frühere Fassung dieses Papers beanspruchte eine solche Abschirmung; dieser Anspruch wird zurückgenommen, und die Verträglichkeit mit lokalen Gravitationstests bleibt offen (Paper~III \cite{Geiger2026c}).
\end{itemize}

Das resultierende Lagrangian $\mathcal{L} = R/(16\pi G) + \gamma R^2 + \frac{1}{2}(\partial\phi)^2 - V_{\mathrm{PT}}(\phi)$ hat seine \textit{Form} vollständig durch die zwei Axiome bestimmt; nur die \textit{numerischen Werte} von $\gamma$, $V_0$ und $\phi_0$ werden durch Daten festgelegt. Diese deduktive Geschlossenheit -- von Annahmen zu einem eindeutigen Lagrangian zu falsifizierbaren Vorhersagen -- erhebt das CRM von einem phänomenologischen Fit zu einem theoretischen Rahmenwerk.

\textit{Reichweite der Serie.} Beansprucht werden die Hintergrund-Expansionsgeschichte und die in den Papers~I--III berichteten kosmologischen Fits. Der $f(R)$-Störungssektor und der lokale Gravitationssektor werden \textit{nicht} beansprucht. Zwei Punkte sind dabei offen zu benennen. Erstens erfüllt der analytische Sektor mit $\gamma \sim H_0^{-2}$, aus dem die Skalaron-Masse und die Compton-Wellenlänge abgeleitet werden, die Cassini-Schranke für den post-Newtonschen Parameter $\gamma_{\mathrm{PPN}}$ nicht: Ohne Abschirmmechanismus bleibt die fünfte Kraft im Sonnensystem ungedämpft (Paper~III \cite{Geiger2026c}). Zweitens enthält keiner der publizierten Fits dieser Serie $\gamma$ als Parameter: Der Planck-Fit aus Paper~III variiert $(\alpha_{M,0}, n)$ neben den kosmologischen Standardparametern, der Pantheon+-Fit aus Paper~II variiert $(\Omega_m, k, a_{\mathrm{trans}}, M)$. Die kosmologischen Fits werden damit von einer eigenständigen effektiven Feldtheorie-Parametrisierung des Störungssektors getragen und nicht von einer numerischen Auswertung des analytischen Sektors; es sind zwei verschiedene Parametrisierungen und nicht zwei Zweige derselben Rechnung. Übersetzt man die Beschränkung $\alpha_{M,0} \sim 10^{-3}$ in die hier verwendete $f(R)$-Identifikation, ergibt sich eine Kopplung, die rund vier Größenordnungen kleiner ist als der analytische Sektor annimmt, und mit entgegengesetztem Vorzeichen. Ob ein einziges $f(R)$-Modell beide Regime abdeckt, ist eine offene Frage dieser Serie und wird hier nicht entschieden.

% ===================================================================
% 9. FAZIT UND AUSBLICK
% ===================================================================
\section{Fazit und Ausblick}
\label{sec:conclusion}

Die vorliegende Arbeit hat gezeigt:

\begin{enumerate}
\item Ein spieltheoretischer Rahmen für die Kosmologie -- das Nash-Gleichgewicht zwischen Nullraum und Raumzeitblase -- führt auf natürliche Weise zu einem Modell, in dem physikalische Gesetze als emergente Gleichgewichtsbedingungen erscheinen. Die spieltheoretischen Axiome besitzen eine exakte thermodynamische Äquivalenz (Nash-Gleichgewicht $\leftrightarrow$ thermodynamisches Gleichgewicht, $\tanh$-Sättigung $\leftrightarrow$ MEPP-Dynamik), die das Rahmenwerk in die Jacobson-Tradition der Gravitation als thermodynamische Zustandsgleichung \cite{Jacobson1995} stellt.
\item Das daraus abgeleitete \textit{Curvature Relaxation Model} (CRM) erklärt die beschleunigte Expansion ohne Dunkle Energie und besteht den Test gegen 1590 Pantheon+ SN-Ia-Datenpunkte nach dem $z>0{,}01$-Schnitt \cite{Scolnic2022}: $\Delta\chi^2 = -12{,}2$ (diagonal) bzw.\ $-11{,}2$ (volle Kovarianzmatrix), $\Delta\mathrm{AIC} = -8{,}2$ bzw.\ $-7{,}2$, vergleichbare Kreuzvalidierung.
\item Die Modellselektion ist gemischt: $\chi^2$ und AIC bevorzugen CRM, BIC bevorzugt $\Lambda$CDM, und die 5-Fold-Kreuzvalidierung ist ohne Fold-Varianzen nicht entscheidend. Vier alternative Funktionalformen liefern $\Delta\chi^2 \approx -9$ bis $-12$ und reduzieren, beseitigen aber nicht das Risiko, dass der Gewinn zusätzliche Flexibilität im späten Expansionsverlauf widerspiegelt.
\item MCMC-basierte Parameterunsicherheiten ($\Omega_m = 0{,}368 \pm 0{,}024$) und die Analyse des homogenen Hintergrunds (kein Big Rip, asymptotisch de-Sitter) stützen die phänomenologische Konsistenz; publikationsreife Konvergenzdiagnostik und ein direkter CPL/CMB/BAO-Vergleich bleiben notwendig.
\item Das CRM macht testbare Vorhersagen: einen späteren Beschleunigungsübergang ($z_{\mathrm{acc}} = 0{,}52$ vs.\ $0{,}84$ in $\Lambda$CDM) und eine charakteristische effektive Zustandsgleichungsgeschichte, die gegen DESI-BAO-, CMB- und Supernova-Kombinationen geprüft werden muss, statt als persistente physische Phantomhaftigkeit interpretiert zu werden. Bereits jetzt findet die CRM-Vorhersage einer verlängerten Wachstumsphase empirische Unterstützung durch JWST-Beobachtungen unerwartet massereicher Galaxien bei hohen Rotverschiebungen \cite{Labbe2023, BoylanKolchin2023} und die statistisch un\-wahrscheinliche Existenz massiver Haufen wie El~Gordo \cite{Asencio2023}.
\item Die Konvergenz dreier unabhängiger Ansätze (CRM, Finsler-Gravitation, Cosmological Teleodynamics) deutet auf einen möglichen Paradigmenwechsel hin: \textit{Dunkle Energie als eigenständige Entität könnte überflüssig sein.}
\end{enumerate}

\textbf{Nächste Schritte} umfassen: (a)~Test gegen Planck-CMB- und DESI-BAO-Daten (die volle Pantheon+-Kovarianzmatrix wurde bereits berücksichtigt), (b)~Berechnung von CMB-Leistungsspektrum und Strukturwachstumsvorhersagen ($f\sigma_8$), (c)~Erforschung der Verbindung zwischen CRM und Finsler-Geometrie, (d)~Entwicklung einer kovarianten Formulierung von $\Phi(a)$ aus dem Ricci-Skalar~$R$, (e)~Untersuchung quantenmechanischer Grundlagen des Krümmungs-Rückgabepotentials, und (f)~Kombination mit lokalen Entfernungsleiter-Daten zur direkten $H_0$-Bestimmung.

\textbf{Ausblick: Vereinigung mit MOND -- Ein Universum ohne dunklen Sektor?} Eine besonders faszinierende Perspektive eröffnet sich durch die Kombination des CRM mit \textit{Modified Newtonian Dynamics} (MOND) \cite{Milgrom1983}. Während das CRM die Dunkle Energie als geometrischen Effekt eliminiert, ersetzt MOND die Dunkle Materie durch eine modifizierte Gravitationsdynamik auf galaktischen Skalen. Beide Rahmenwerke konvergieren in der Vorhersage, dass Strukturen früher und effizienter entstehen als in $\Lambda$CDM -- das CRM durch eine verlängerte materiedominierte Ära ($z_{\mathrm{acc}} = 0{,}52$), MOND durch effektiv stärkere Gravitation bei niedrigen Beschleunigungen \cite{Asencio2023}. Eine vorläufige Analyse mit einem rein baryonischen Universum ($\Omega_m = \Omega_b \approx 0{,}05$) und einem erweiterten geometrischen Potential $\Omega_\Phi(a) = \Phi_0 \cdot f_{\mathrm{tanh}}(a) + \alpha \cdot a^{-\beta}$ liefert $\Delta\chi^2 = -26{,}3$ und $\Delta\mathrm{AIC} = -16{,}3$ gegenüber $\Lambda$CDM -- \textit{deutlich besser als sowohl das Standard-CRM als auch $\Lambda$CDM}. Die MCMC-Posterioranalyse ergibt $\beta = 2{,}02 \pm 0{,}20$, was exakt der Skalierung räumlicher Krümmung ($a^{-2}$) entspricht. Die ``Dunkle Materie'' wäre demnach ein dynamischer Krümmungseffekt, kein Teilchen. Dieses Ergebnis bedarf einer vollständigen relativistischen Behandlung (z.\,B.\ TeVeS oder AeST \cite{Bekenstein2004,Skordis2021}) und wird in einer Folgearbeit detailliert analysiert.

\begin{quote}
\textit{``Manchmal ist die eleganteste Erklärung nicht eine neue Kraft, sondern eine nachlassende Einschränkung.''}
\end{quote}

% ===================================================================
% LITERATUR
% ===================================================================
\subsection*{Software}

Diese Arbeit verwendet \texttt{emcee} \cite{ForemanMackey2013} für MCMC-Sampling, \texttt{NumPy} \cite{Harris2020}, \texttt{SciPy} \cite{Virtanen2020} für numerische Berechnungen und \texttt{Matplotlib} \cite{Hunter2007} für Visualisierungen; \texttt{CLASS} \cite{Blas2011} dient als Referenzcode für die geplanten CMB-Vergleiche.

\begin{acknowledgments}
Der Autor dankt den Entwicklern von \texttt{CLASS} \cite{Blas2011} und
\texttt{emcee} \cite{ForemanMackey2013} für die Bereitstellung ihrer Codes
als Open Source.

\paragraph{KI-Werkzeuge.}
KI-basierte Werkzeuge (Claude von Anthropic; Copilot von Microsoft/GitHub; Gemini von Google DeepMind)
wurden für Codeentwicklung, numerische Analysen, Literaturrecherche und redaktionelle
Unterstützung eingesetzt. Alle physikalischen Hypothesen, mathematischen Ableitungen und die wissenschaftliche
Interpretation sind ausschliesslich das Werk des Autors.
\end{acknowledgments}

\paragraph{Finanzierung.}
Diese Forschung erhielt keine externe Finanzierung.

\paragraph{Datenverfügbarkeit.}
Die Pantheon+ Supernova-Daten \cite{Scolnic2022} sind öffentlich zugänglich.
Der gesamte Analyse-Code sowie alle numerischen Ergebnisse stehen unter
\url{https://github.com/research-line/crm-cosmology}
(DOI: \href{https://doi.org/10.5281/zenodo.18728936}{10.5281/zenodo.18728936}) zur Verfügung.

\begin{thebibliography}{99}

\bibitem{Riess1998}
Riess, A.\,G.\ et al.\ (1998).
Observational Evidence from Supernovae for an Accelerating Universe and a Cosmological Constant.
\textit{The Astronomical Journal}, 116(3), 1009--1038.
DOI: 10.1086/300499.

\bibitem{Perlmutter1999}
Perlmutter, S.\ et al.\ (1999).
Measurements of $\Omega$ and $\Lambda$ from 42 High-Redshift Supernovae.
\textit{The Astrophysical Journal}, 517(2), 565--586.
DOI: 10.1086/307221.

\bibitem{Planck2020}
Planck Collaboration (2020).
Planck 2018 results. VI. Cosmological parameters.
\textit{Astronomy \& Astrophysics}, 641, A6.
DOI: 10.1051/0004-6361/201833910.

\bibitem{Weinberg1989}
Weinberg, S.\ (1989).
The Cosmological Constant Problem.
\textit{Reviews of Modern Physics}, 61(1), 1--23.
DOI: 10.1103/RevModPhys.61.1.

\bibitem{Riess2022}
Riess, A.\,G.\ et al.\ (2022).
A Comprehensive Measurement of the Local Value of the Hubble Constant with 1\,km/s/Mpc Uncertainty from the Hubble Space Telescope and the SH0ES Team.
\textit{The Astrophysical Journal Letters}, 934(1), L7.
DOI: 10.3847/2041-8213/ac5c5b.

\bibitem{DESI2024}
DESI Collaboration (2025).
DESI 2024 VI: Cosmological Constraints from the Measurements of Baryon Acoustic Oscillations.
\textit{Journal of Cosmology and Astroparticle Physics}, 2025(02), 021.
DOI: 10.1088/1475-7516/2025/02/021.

\bibitem{Pfeifer2025}
Pfeifer, C.\ et al.\ (2025).
From kinetic gases to an exponentially expanding universe -- the Finsler-Friedmann equation.
\textit{Journal of Cosmology and Astroparticle Physics}, 2025(10), 050.
DOI: 10.1088/1475-7516/2025/10/050.

\bibitem{Trivedi2025}
Trivedi, O.\ \& Venkatasubramanian, V.\ (2026).
Game theory in cosmology.
\textit{Physics of the Dark Universe}, 52, 102277.
DOI: 10.1016/j.dark.2026.102277.

\bibitem{Caldwell1998}
Caldwell, R.\,R., Dave, R.\ \& Steinhardt, P.\,J.\ (1998).
Cosmological Imprint of an Energy Component with General Equation of State.
\textit{Physical Review Letters}, 80(8), 1582--1585.
DOI: 10.1103/PhysRevLett.80.1582.

\bibitem{Starobinsky1980}
Starobinsky, A.\,A.\ (1980).
A New Type of Isotropic Cosmological Models Without Singularity.
\textit{Physics Letters B}, 91(1), 99--102.
DOI: 10.1016/0370-2693(80)90670-X.

\bibitem{Sotiriou2010}
Sotiriou, T.\,P.\ \& Faraoni, V.\ (2010).
$f(R)$ Theories of Gravity.
\textit{Reviews of Modern Physics}, 82(1), 451--497.
DOI: 10.1103/RevModPhys.82.451.

\bibitem{Verlinde2011}
Verlinde, E.\ (2011).
On the Origin of Gravity and the Laws of Newton.
\textit{Journal of High Energy Physics}, 2011, 29.
DOI: 10.1007/JHEP04(2011)029.

\bibitem{Verlinde2017}
Verlinde, E.\ (2017).
Emergent Gravity and the Dark Universe.
\textit{SciPost Physics}, 2(3), 016.
DOI: 10.21468/SciPostPhys.2.3.016.

\bibitem{Padmanabhan2012}
Padmanabhan, T.\ (2012).
Emergent Perspective of Gravity and Dark Energy.
\textit{Research in Astronomy and Astrophysics}, 12(8), 891.
DOI: 10.1088/1674-4527/12/8/003.

\bibitem{Euclid2024}
Euclid Collaboration (2025).
Euclid Quick Data Release (Q1) -- Data release overview.
ESA/Euclid Consortium.
DOI: 10.48550/arXiv.2503.15302.

\bibitem{Casimir1948}
Casimir, H.\,B.\,G.\ (1948).
On the attraction between two perfectly conducting plates.
\textit{Proceedings of the Royal Netherlands Academy of Arts and Sciences}, 51, 793--795.

\bibitem{Hawking1974}
Hawking, S.\,W.\ (1974).
Black hole explosions?
\textit{Nature}, 248, 30--31.
DOI: 10.1038/248030a0.

\bibitem{Nash1950}
Nash, J.\,F.\ (1950).
Equilibrium points in $n$-person games.
\textit{Proceedings of the National Academy of Sciences}, 36(1), 48--49.
DOI: 10.1073/pnas.36.1.48.

\bibitem{DESI2025}
DESI Collaboration (2025).
DESI DR2 Results II: Measurements of Baryon Acoustic Oscillations and Cosmological Constraints.
\textit{Physical Review D}, 112, 083515.
DOI: 10.1103/tr6y-kpc6.

\bibitem{Efstratiou2026}
Efstratiou, D., Paraskevas, E.\,A.\ \& Perivolaropoulos, L.\ (2026).
Addressing the DESI DR2 phantom-crossing anomaly and enhanced $H_0$ tension with reconstructed scalar-tensor gravity.
\textit{Physical Review D}, 113(12), 123525.
DOI: 10.1103/PhysRevD.113.123525.

\bibitem{Scolnic2022}
Scolnic, D.\ et al.\ (2022).
The Pantheon+ Analysis: The Full Data Set and Light-curve Release.
\textit{The Astrophysical Journal}, 938(2), 113.
DOI: 10.3847/1538-4357/ac8b7a.

\bibitem{KassRaftery1995}
Kass, R.\,E.\ \& Raftery, A.\,E.\ (1995).
Bayes Factors.
\textit{Journal of the American Statistical Association}, 90(430), 773--795.
DOI: 10.1080/01621459.1995.10476572.

\bibitem{ForemanMackey2013}
Foreman-Mackey, D.\ et al.\ (2013).
emcee: The MCMC Hammer.
\textit{Publications of the Astronomical Society of the Pacific}, 125(925), 306--312.
DOI: 10.1086/670067.

\bibitem{Labbe2023}
Labb\'e, I.\ et al.\ (2023).
A population of red candidate massive galaxies $\sim$600\,Myr after the Big Bang.
\textit{Nature}, 616(7956), 266--269.
DOI: 10.1038/s41586-023-05786-2.

\bibitem{BoylanKolchin2023}
Boylan-Kolchin, M.\ (2023).
Stress testing $\Lambda$CDM with high-redshift galaxy candidates.
\textit{Nature Astronomy}, 7, 731--735.
DOI: 10.1038/s41550-023-01937-7.

\bibitem{Asencio2023}
Asencio, E., Banik, I.\ \& Kroupa, P.\ (2023).
The El Gordo galaxy cluster challenges $\Lambda$CDM for any plausible collision velocity.
\textit{The Astrophysical Journal}, 954(2), 162.
DOI: 10.3847/1538-4357/ace62a.

\bibitem{Miller2018}
Miller, T.\,B.\ et al.\ (2018).
A massive core for a cluster of galaxies at a redshift of 4.3.
\textit{Nature}, 556(7702), 469--472.
DOI: 10.1038/s41586-018-0025-2.

\bibitem{Milgrom1983}
Milgrom, M.\ (1983).
A modification of the Newtonian dynamics as a possible alternative to the hidden mass hypothesis.
\textit{The Astrophysical Journal}, 270, 365--370.
DOI: 10.1086/161130.

\bibitem{Skordis2021}
Skordis, C.\ \& Z{\l}o\'snik, T.\ (2021).
New Relativistic Theory for Modified Newtonian Dynamics.
\textit{Physical Review Letters}, 127(16), 161302.
DOI: 10.1103/PhysRevLett.127.161302.

\bibitem{Bekenstein2004}
Bekenstein, J.\,D.\ (2004).
Relativistic gravitation theory for the modified Newtonian dynamics paradigm.
\textit{Physical Review D}, 70(8), 083509.
DOI: 10.1103/PhysRevD.70.083509.

\bibitem{HuSawicki2007}
Hu, W.\ \& Sawicki, I.\ (2007).
Models of $f(R)$ Cosmic Acceleration that Evade Solar-System Tests.
\textit{Physical Review D}, 76(6), 064004.
DOI: 10.1103/PhysRevD.76.064004.

\bibitem{Jacobson1995}
Jacobson, T.\ (1995).
Thermodynamics of Spacetime: The Einstein Equation of State.
\textit{Physical Review Letters}, 75(7), 1260--1263.
DOI: 10.1103/PhysRevLett.75.1260.

\bibitem{JacobsonNonRiemannian2026}
Martinez, J.\,N., Rodriguez-Ruiz, J.\,F.\ \& Rodriguez, Y.\ (2026).
Jacobson's thermodynamic approach to classical gravity applied to non-Riemannian geometries: remarks on the simplicity of Nature.
\textit{Preprint}, arXiv:2602.00422.

\bibitem{Dewar2003}
Dewar, R.\ (2003).
Information theory explanation of the fluctuation theorem, maximum entropy production and self-organized criticality in non-equilibrium stationary states.
\textit{Journal of Physics A: Mathematical and General}, 36(3), 631--641.
DOI: 10.1088/0305-4470/36/3/303.

\bibitem{Wolpert2006}
Wolpert, D.\,H.\ (2006).
Information Theory -- The Bridge Connecting Bounded Rational Game Theory and Statistical Physics.
In: \textit{Complex Engineered Systems}, Springer, 262--290.
DOI: 10.1007/3-540-32834-3\_12.

\bibitem{Harris2020}
Harris, C.\,R.\ et al.\ (2020).
Array programming with NumPy.
\textit{Nature}, 585, 357--362.
DOI: 10.1038/s41586-020-2649-2.

\bibitem{Virtanen2020}
Virtanen, P.\ et al.\ (2020).
SciPy 1.0: Fundamental Algorithms for Scientific Computing in Python.
\textit{Nature Methods}, 17, 261--272.
DOI: 10.1038/s41592-019-0686-2.

\bibitem{Hunter2007}
Hunter, J.\,D.\ (2007).
Matplotlib: A 2D Graphics Environment.
\textit{Computing in Science \& Engineering}, 9(3), 90--95.
DOI: 10.1109/MCSE.2007.55.

\bibitem{GelmanRubin1992}
Gelman, A.\ \& Rubin, D.~B.\ (1992).
Inference from Iterative Simulation Using Multiple Sequences.
\textit{Statistical Science}, 7(4), 457--472.
DOI: 10.1214/ss/1177011136.

\bibitem{Chevallier2001}
Chevallier, M.\ \& Polarski, D.\ (2001).
Accelerating Universes with Scaling Dark Matter.
\textit{Int.\ J.\ Mod.\ Phys.\ D}, 10, 213.
DOI: 10.1142/S0218271801000822.

\bibitem{Linder2003}
Linder, E.~V.\ (2003).
Exploring the Expansion History of the Universe.
\textit{Phys.\ Rev.\ Lett.}, 90, 091301.
DOI: 10.1103/PhysRevLett.90.091301.

\bibitem{Blas2011}
Blas, D., Lesgourgues, J.\ \& Tram, T.\ (2011).
The Cosmic Linear Anisotropy Solving System (CLASS) II: Approximation schemes.
\textit{Journal of Cosmology and Astroparticle Physics}, 2011(07), 034.
DOI: 10.1088/1475-7516/2011/07/034.

\bibitem{Geiger2026c}
Geiger, L.\ (2026).
Microscopic Foundations of the Curvature Relaxation Model: From Quantum Geometry to Macroscopic Saturation.
Companion paper (Paper~III), submitted simultaneously.

\bibitem{Geiger2026d}
Geiger, L.\ (2026).
The Galactic--Cosmological Nexus: Connecting CRM to MOND Dynamics via Curvature Saturation.
Companion paper (Paper~IV), submitted simultaneously.

\bibitem{Kovacs2026}
Kovacs, G.\ (2026).
Gravity Imprints, Time Fades: A Curvature-Memory Relaxation Model with Exact MOND Steady State.
Non-peer-reviewed manuscript v4.4, March 2026.
Available at \url{https://vixra.org/abs/2604.0006}.

\bibitem{Stichel2016}
Stichel, P.\,C.\ (2016).
Cosmological model with dynamical curvature.
\textit{Preprint}, arXiv:1601.07030v3 [physics.gen-ph].
DOI: 10.48550/arXiv.1601.07030.

\bibitem{EstebanGutierrez2025}
Esteban-Guti\'errez, A., Garc\'ia-Aspeitia, M.\,A., Hern\'andez-Almada, A., Maga\~na, J.\ \& Motta, V.\ (2025).
On a model of variable curvature that mimics the observed Universe acceleration.
\textit{Preprint}, arXiv:2410.08306v2 [astro-ph.CO].
DOI: 10.48550/arXiv.2410.08306.

\end{thebibliography}

\end{document}
